\documentclass[11pt,a4paper]{article}

\usepackage[T1]{fontenc}
\usepackage[utf8]{inputenc}
\usepackage[british]{babel}
\usepackage{lmodern}
\usepackage{microtype}
\usepackage[a4paper,margin=1in]{geometry}
\usepackage{amsmath,amssymb,amsthm,mathtools}
\usepackage{enumitem}
\usepackage{booktabs,tabularx,array}
\usepackage{xcolor}
\usepackage{hyperref}
\hypersetup{
  colorlinks=true,
  linkcolor=black,
  citecolor=black,
  urlcolor=black,
  hypertexnames=false,
  pdftitle={A Preference for Traceable Agency: An Axiomatization of Expected Utility plus Mutual Information},
  pdfauthor={Daniele Condorelli},
  pdfsubject={A behavioural characterisation of expected utility plus mutual information}
}
\setlength{\emergencystretch}{2em}
\interfootnotelinepenalty=10000

\newtheorem{theorem}{Theorem}
\newtheorem{proposition}[theorem]{Proposition}
\newtheorem{lemma}[theorem]{Lemma}
\newtheorem{corollary}[theorem]{Corollary}
\theoremstyle{definition}
\newtheorem{definition}[theorem]{Definition}
\newtheorem{remark}[theorem]{Remark}

\newcommand{\D}{\Delta}
\newcommand{\E}{\mathbb E}
\newcommand{\I}{I}
\newcommand{\Sh}{H}
\newcommand{\supp}{\operatorname{supp}}
\newcommand{\Id}{\operatorname{Id}}
\newcommand{\Rnum}{\mathbb R}
\newcommand{\R}{\mathbb R}
\newcommand{\KL}{D_{\mathrm{KL}}}
\newcommand{\1}{\mathbf 1}

\title{A Preference for Traceable Agency\\[0.35em]
\large An Axiomatization of Expected Utility plus Mutual Information\thanks{Email:
\href{mailto:d.condorelli@gmail.com}{d.condorelli@gmail.com}.  I thank the LSE Theory Group for hospitality and stimulating discussion during the summer of 2025, when this project was first conceived. I thank Peter Hammond and Agustin Troccoli Moretti for comments. Joe Basford suggested a way to simplify Axiom~(A8).

\smallskip
All parts of this manuscript have been written with assistance from a multi-model LLM workflow.  The
question, the model, the axioms, and the form of the representation are the work of human hands.  The author
personally contributed very little to the writing of the proofs beyond advice and encouragement.  An earlier
machine-checked Lean~4 proof of Theorem~\ref{thm:main} is available at
\url{https://github.com/dcond/TraceableAgency}: the verified declaration is
\href{https://github.com/dcond/TraceableAgency/blob/main/TraceableAgency/Theorem1/Proof.lean}{\texttt{trace\_tempered\_choice\_v10\_theorem1}},
and the \href{https://github.com/dcond/TraceableAgency/blob/main/Certificate/FORMALIZATION_CERTIFICATE.md}{formalisation certificate}
maps the axioms and theorem to their formal statements.  The shorter proof presented here, together with the
uniqueness, sign, and independence results in Section~\ref{sec:theorem}, has also been machine-checked in
Lean~4.  All claims, errors, and omissions remain the sole responsibility of the author.}}
\author{Daniele Condorelli\\University of Warwick}
\date{\today}

\begin{document}
\maketitle

\begin{abstract}
\noindent
An individual may value her actions both for their payoffs and for what their observable consequences reveal
about which one she took.  A channel maps actions to distributions over consequences; in each channel she ranks
lotteries over actions.  Three axioms are distinctive: she never prefers her action garbled, nor its consequence, and
a sure-thing principle holds with interim records as events.  Together with standard ones, they characterise
expected payoff utility plus a positive, channel-independent weight on the mutual information between action and
consequence, her trace.  The consequences of a predetermined act reveal nothing new, so randomising can be
strictly optimal.
\end{abstract}

\medskip
\noindent\textit{Keywords:} traceable agency; mutual information; Shannon entropy; preference for
randomisation; intrinsic preference for information; non-expected utility.\\
\textit{JEL codes:} D81, D83, D91.

\clearpage

% ===============================================================
\section{Introduction}
% ===============================================================


\begin{quote}
\emph{Shower upon him every earthly blessing \ldots\ He would even risk his cakes and would deliberately desire the
most fatal rubbish, the most uneconomical absurdity \ldots\ simply in order to prove to himself---as though that were
so necessary---that men still are men and not the keys of a piano.\ \ldots\ If you say that all this, too, can be
calculated and tabulated \ldots\ then man would purposely go mad in order to be rid of reason and gain his point!}

\hfill --- Fyodor Dostoevsky, \emph{Notes from the Underground}, Part I, Chapter VIII \cite{Dostoevsky1864}
\end{quote}


Dostoevsky's Underground Man faces a choice problem.  The cakes are material payoff, and he will risk them so
that his conduct can attest to being his own, not the mechanical implication of a table.  The proof he seeks
is self-defeating, because an act fixed in advance cannot provide it and deliberate rebellion can be
“calculated and tabulated.”  So he would “purposely go mad in order to be rid of reason”---not because
madness is what he values, but because it is his last imagined way of preventing the table from closing before he
acts.  This paper reads madness as making oneself unpredictable, and takes the evidentiary idea literally.  In
Dostoevsky's example conduct is observed directly; we consider the general case, in which an act is observed
only through its visible consequences.

For visible consequences to carry information about an act, two conditions are needed.  There must be uncertainty
beforehand about which act will occur, and the consequence must discriminate among the
possible acts.  A lottery over acts leaves no trace when every act generates the same distribution of
consequences; a consequence that perfectly reveals the act conveys nothing new when the act was already
known.  When both conditions hold, the consequence is evidence about the act.  We call an intrinsic
preference for such evidence a preference for \emph{traceable agency}.  The paper axiomatises it jointly
with ordinary preferences over payoff risk.

Let $A$ be a finite set of actions, $O$ a finite set of payoff-bearing
outcomes, $R$ a finite set of explicit visible records, and $K:A\to\Delta(O\times R)$ the channel relating actions to
visible consequences.  An action lottery $q\in\Delta(A)$ gives the distribution of the realised action.  The primitive is a family
$\{\succeq_K\}_K$, with one binary relation $\succeq_K$ on $\Delta(A)$ for each fixed channel $K$.
The statement $q\succeq_Kp$ means that $q$ is weakly preferred to $p$ in $K$.  Eight axioms are
imposed on this family.  To compare lotteries governed
by different channels, the channels are placed side by side as blocks of one larger environment; duplicating a
channel or adding an unused block does not alter a comparison.  Two processing axioms say that the agent never gains when her record is garbled or her action variable is processed noisily, and a benchmark comparison rules out trace indifference. A recordwise sure-thing principle for compounded records supplies
the cardinal structure.  The taste prices how much the consequence reveals, not what it reveals.

A family $\{\succeq_K\}_K$ satisfies the eight axioms if and
only if there are a nonconstant payoff utility $u:O\to\mathbb R$ and $\lambda>0$ such that, for every finite joint
channel $K$ and every $q,p\in\Delta(A)$,
\[
 q\succeq_Kp
 \quad\Longleftrightarrow\quad
 \E_{qK}[u(O)]+\lambda I_{qK}(A;O,R)
 \ge
 \E_{pK}[u(O)]+\lambda I_{pK}(A;O,R),
\]
where $I_{qK}(A;O,R)$ is the mutual information between the action and its visible consequence.  We call such
families \emph{trace-tempered}.  The utility $u$ is unique up to positive affine transformations; after fixing that
scale and the logarithm base, $\lambda$ is measured in payoff units per unit of information, and the same pair prices comparisons across channels.  The sign is set by three axioms: reversing them characterises $\lambda<0$, a preference for statistical
deniability, and replacing them by indifference clauses yields expected utility, $\lambda=0$.

While there is no observer in the model, preferences admit an instrumental reading: an audience rewarding the agent in proportion to what the visible consequence lets it learn about the act induces the criterion. Rather than a taste, $\lambda$ becomes a price of attribution, positive for a politician backed when acts move outcomes, negative if what is learned is used against~her.


The information term has a simple interpretation.  Observing $(O,R)$ induces a posterior belief about the realised
action; traceability is the expected reduction in that uncertainty, measured by Shannon entropy $\Sh$.  The value is
zero if visible consequences are independent of actions, and zero again if the action is known in advance
because $q$ is degenerate.  If the visible-consequence supports of the actions in the support of
$q$ are pairwise disjoint, the value is $\Sh(q)$---Dostoevsky's case, in which conduct is observed directly.\footnote{If the
visible-consequence supports of all channel rows are pairwise disjoint, the unique optimal lottery has
multinomial-logit weights.  Overlap can make independence of irrelevant alternatives fail; see
\S\ref{sec:randomisation}.}  Thus the criterion is
not a taste for randomisation as such.  Randomisation matters only because it creates uncertainty about the act for the
visible consequence to resolve.

Although the axioms impose further structure on how trace is valued,
the underlying taste is for the coupling of act and consequence.  A
four-action channel turns this into a direct test.  Every action pays the same amount and writes one of two
symbols on the record.  Action $a_1$ always writes the first symbol and $a_2$ the second; $a_3$ and $a_4$ each
write either symbol with equal probability.  The choice is between a fair coin over $\{a_1,a_2\}$ and a fair
coin over $\{a_3,a_4\}$.  Both coins are fair and induce the same joint distribution of payoff
and record.  Under the first coin the record identifies the realised action; under the second it reveals
nothing.  Any criterion that reads only the distribution of consequences is indifferent between the coins; so is
a taste for randomisation as such.  Traceable agency predicts a strict preference for the first.  Because the
record affects neither payoffs nor any later decision, the preference cannot be
instrumental.\footnote{A second treatment completes the design: repeat
the choice in a channel that writes a fair record independently of the realised action.  Every representation
then assigns both flips the same value, whatever $u$ and $\lambda$, so a strict preference there rejects the
model.  Section~\ref{sec:implications} develops the test and what it excludes.}  The comparison identifies a
value of act--consequence dependence; by itself it does not distinguish trace from control or identify the
psychological source of that value.\footnote{The information term admits a control interpretation. Read backward, it measures trace. Read forward, it measures how strongly the action controls the distribution of visible consequences. We return to this in Section~\ref{sec:interpretation}.}


This comparison has not been run.  The nearest evidence comes from two experiments.  Bartling, Fehr, and Herz elicit each principal's point of indifference between retaining and delegating the right to choose a project and the effort devoted to it.  At that point, the delegation lottery has a certainty equivalent 16.7\% higher on average than the lottery produced by the principal's own choices, when both are valued outside the delegation setting \cite{BartlingFehrHerz2014}.  Principals pay to keep the payoff lottery dependent on their own choices, which a trace reading accommodates only if that choice is itself uncertain. The same evidence may instead reflect autonomy.  Mistry and Liljeholm, and later Norton and Liljeholm, ask subjects to select a room containing two slot machines and then gamble in that room.  Subjects are more likely to choose a high-divergence room over a zero-divergence room when they choose between its machines themselves than when a computer chooses for them, holding current expected payoffs and outcome entropies fixed \cite{MistryLiljeholm2016,NortonLiljeholm2020}.  With equal weight on the two machines, the divergence between their token distributions is the mutual information between the machine and the token.  This also admits a trace interpretation: in self-play, the token is evidence about a voluntarily selected machine.  But token values could change after the room was selected, allowing subjects in self-play to adapt their machine choices.  Thus trace is confounded with the value of deciding in the first design and with that of adapting in the second.  The test proposed above removes both confounds.

Section~\ref{sec:model} sets out the domain and the eight axioms.  Section~\ref{sec:theorem} states the representation theorem,
sketches the proof, and derives uniqueness, the sign trichotomy, and the
independence of the axioms.
Section~\ref{sec:implications} presents a two-treatment test that can reject the model; derives deliberate randomisation, with multinomial logit as the boundary case when channel-row supports are pairwise disjoint; measures $\lambda$ from choices; and applies the criterion to Holmstr\"om's moral-hazard problem.
Section~\ref{sec:literature} reviews the literature.
Section~\ref{sec:interpretation} discusses interpretation and scope.  The proofs are in the appendices:
Appendix~\ref{app:main-proof} gives a proof of Theorem~1, and Appendix~\ref{app:auxiliary-results} contains the remaining results.


% ===============================================================
\section{The model}\label{sec:model}
% ===============================================================

Fix a finite set $O$ of \emph{payoff-bearing outcomes}, with $|O|\ge2$.  Let $A$ be a nonempty finite set of
\emph{actions} and $R$ a nonempty finite set of \emph{explicit visible records}.  A joint channel is a stochastic
kernel
\[
  K:A\longrightarrow\D(O\times R).
\]
For each finite joint channel $K$, the primitive behavioural datum is a binary relation
\[
 \succeq_K\quad\text{on}\quad\D(A).
\]
The interpretation is that $q$ is a lottery over actions used by the agent and $q\succeq_Kp$ means that, in the fixed
environment $K$, $q$ is evaluated at least as highly as $p$.  Within that environment, the lottery is the
procedure being evaluated and the channel is fixed.\footnote{Why not a single
preference over joint laws on $A\times O\times R$?  Every joint law factors as $q(a)K(o,r\mid a)$ for some pair,
so the domain is rich enough.  But in any one choice situation the channel is fixed: the agent moves the
marginal on $A$, never the coupling.  A single preference would rank couplings, and no choice reveals those
rankings.}

\bigskip
Axioms are imposed on the family $\{\succeq_K\}_K$.
The first two axioms are standard, and Axiom~\textup{(A3)} is the usual nondegeneracy condition on material
outcomes.  Axioms~\textup{(A4)}, \textup{(A6)}, and \textup{(A7)} carry the preference for trace.
Axiom~\textup{(A5)} puts the separate within-channel rankings on a common scale, while
Axiom~\textup{(A8)} imposes consistency under compounding; neither fixes the sign of the trace term.

\begin{description}[leftmargin=*,style=sameline]

\item[(A1) Weak order.]
For every finite joint channel $K$, $\succeq_K$ is complete and transitive on $\D(A)$.

\item[(A2) Continuity.]
For each fixed triple of alphabets $(A,O,R)$, the set
\[
 \{(K,q,p):q\succeq_Kp\}
 \subseteq \D(O\times R)^A\times\D(A)^2
\]
is closed.


\item[(A3) Material relevance.]
There are distinct outcomes $o^+,o^-\in O$ and a two-action channel $K$ with a one-point record set in which $a^+$ delivers
$o^+$ for sure and $a^-$ delivers $o^-$ for sure.  For this channel,
\[
 \delta_{a^+}\succ_K\delta_{a^-}.
\]
The agent strictly ranks at least two outcomes when each is delivered for sure and the two actions leave the same record.

\item[(A4) Trace relevance.]
There is an outcome $o_\ast\in O$ and a four-action channel $K$ in which every action delivers $o_\ast$ for
sure.  Actions $a_1$ and $a_2$ leave distinct records $r_1$ and $r_2$, respectively, while each of $a_3$ and
$a_4$ leaves either record with equal probability.  For this channel,
\[
 \frac12(\delta_{a_1}+\delta_{a_2})
 \succ_K
 \frac12(\delta_{a_3}+\delta_{a_4}).
\]
Both are fair lotteries and produce the same distribution of visible consequences. The axiom requires a strict
preference for the first, whose record identifies the realised action, over the second, whose record does not.

\end{description}

\paragraph{Comparison environments.}
Some axioms compare lotteries carried by different channels.  Each relation $\succeq_K$ ranks lotteries within
one channel, so such comparisons need a domain in which two channels appear in a single preference statement.  A
larger channel that contains both as labelled blocks provides it.

For
$K:A\to\D(O\times R)$ and $L:B\to\D(O\times R_L)$, let $K\sqcup L$ be the block-diagonal joint channel---the
\emph{block environment}---with action set
$(A\times\{0\})\cup(B\times\{1\})$, payoff set $O$, and record set
$(R\times\{0\})\cup(R_L\times\{1\})$: each block operates its own channel, and all cross-block probabilities
are zero.  For $q\in\D(A)$ and $p\in\D(B)$,
let $q^0$ and $p^1$ denote their
block-supported copies, defined on the tagged action set by
\[
\begin{aligned}
 q^0(a,0)&:=q(a), & q^0(b,1)&:=0,\\
 p^1(a,0)&:=0,    & p^1(b,1)&:=p(b),
\end{aligned}
\qquad(a\in A,\ b\in B).
\]

Under \(q^0\), the visible consequence is generated by \(K\) after the action is drawn according to \(q\); under \(p^1\), it is generated by \(L\) after the action is drawn according to \(p\). So a choice between them is
also a choice between the two channels, made inside a single environment.
Write
\[
 (q,K)\succeq(p,L)
 \quad\Longleftrightarrow\quad
 q^0\succeq_{K\sqcup L}p^1.
\]
Each instance of this shorthand is a single comparison inside the two-block environment $K\sqcup L$, which is
built from the two channels.  Here $\succ$ and $\sim$ denote strict preference and indifference between these
two lotteries within that same block environment.  The pair $(q,K)$
has no meaning in isolation.\footnote{Formally, $(K\sqcup L)(o,(r,0)\mid(a,0))=K(o,r\mid a)$ and
$(K\sqcup L)(o,(r_L,1)\mid(b,1))=L(o,r_L\mid b)$, with every other entry zero.  In the binary case, with
$A=\{a_1,a_2\}$ and $B=\{b_1,b_2\}$, the block environment has the four tagged actions $(a_1,0)$, $(a_2,0)$,
$(b_1,1)$, $(b_2,1)$; a lottery $q=(q(a_1),q(a_2))$ becomes $q^0=(q(a_1),q(a_2);0,0)$ and $p=(p(b_1),p(b_2))$
becomes $p^1=(0,0;p(b_1),p(b_2))$, the semicolon separating the blocks.  Then $(q,K)\succeq(p,L)$ abbreviates
$q^0\succeq_{K\sqcup L}p^1$: one ordinary comparison of two lotteries in this one four-action environment.}  The construction extends to any finite family
$\{K_i:A_i\to\D(O\times R_i)\}_{i\in I}$: write $\bigsqcup_{i\in I}K_i$ for the analogous block environment and
$q_i^i$ for the lottery that assigns $q_i(a)$ to $(a,i)$ and zero to every other block.

\begin{description}[leftmargin=*,style=sameline]

\item[(A5) Block-comparison coherence.]
\emph{Duplication:} for every $K:A\to\D(O\times R)$ and all $q,p\in\D(A)$,
\[
 q\succeq_Kp
 \quad\Longleftrightarrow\quad
 (q,K)\succeq(p,K).
\]
\emph{Irrelevant blocks:} for every finite block environment $\bigsqcup_{k\in I}K_k$, every ordered pair of
distinct blocks $i,j\in I$, and all $q_i\in\D(A_i)$ and $q_j\in\D(A_j)$,
\[
 q_i^i\succeq_{\bigsqcup_{k\in I}K_k}q_j^j
 \quad\Longleftrightarrow\quad
 (q_i,K_i)\succeq(q_j,K_j).
\]

Comparing $q$ and $p$ across two tagged copies of $K$ is just comparing them in $K$; adding blocks not involved
in a comparison leaves it unchanged.

\end{description}

\paragraph{Processing.}
The next two axioms use two ways of transforming an environment.  A \emph{record processor} is a stochastic
kernel $T:O\times R\to\D(O\times R')$ that garbles the record without changing the payoff:
$T(o',r'\mid o,r)=0$ whenever $o'\neq o$.  The garbled channel is simply the matrix product
$KT:A\to\D(O\times R')$,
\begin{equation}
 (KT)(o',r'\mid a)
 :=\sum_{o,r}K(o,r\mid a)T(o',r'\mid o,r).
 \label{eq:record-processing}
\end{equation}

An \emph{action processor} is a stochastic kernel $S:A\to\D(B)$ that replaces the action variable $A$ by a
processed action variable $B$.  Given $q\in\D(A)$, it induces the lottery $qS\in\D(B)$,
\[
 (qS)(b):=\sum_aq(a)S(b\mid a).
\]
A \emph{Bayesian pushforward} of $K$ by $S$ at $q$ is any channel $\widehat K:B\to\D(O\times R)$
satisfying\footnote{If $(qS)(b)=0$, the
equation leaves that row unrestricted; every statement below applies to every consistent completion, making
unreached processed actions irrelevant.}
\begin{equation}
 (qS)(b)\,\widehat K(o,r\mid b)
 =\sum_a q(a)S(b\mid a)K(o,r\mid a)
 \quad\text{for every }(b,o,r).
 \label{eq:action-processing}
\end{equation}
In the processed environment, the alternatives are the noisy names $b$, while the visible consequences remain
$(O,R)$.  Choosing $qS$ under $\widehat K$ replaces dependence on the original action with dependence on its
noisy name, without changing the marginal distribution of visible consequences.

\begin{description}[leftmargin=*,style=sameline]

\item[(A6) Record data processing.]
For every $K:A\to\D(O\times R)$, every record processor $T$, and every $q\in\D(A)$,
\[
 (q,K)\succeq(q,KT).
\]
The agent never gains when her record is garbled.

\item[(A7) Action data processing.]
For every $K:A\to\D(O\times R)$, every action processor $S:A\to\D(B)$, every $q\in\D(A)$, and every Bayesian
pushforward $\widehat K$ of $K$ by $S$ at $q$,
\[
 (q,K)\succeq(qS,\widehat K).
\]
Nor does she gain when the action variable is processed, holding the distribution of visible consequences fixed. Special cases include swapping action labels, adding or deleting unused actions, splitting an action into identical copies, merging actions, and noisy recoding.
\end{description}

\paragraph{Compounding.}
For $k\ge1$, let $P:A\to\D(\{1,\ldots,k\})$ be a first-stage record channel and let
$(K_1,\ldots,K_k)$ be an ordered list of continuation channels $K_i:A\to\D(O\times R)$.  Their
\emph{compound environment} is
\begin{equation}
 \bigl(P\triangleright(K_1,\ldots,K_k)\bigr)(o,(i,r)\mid a)
 :=P(i\mid a)K_i(o,r\mid a),
 \qquad i=1,\ldots,k.
 \label{eq:compound}
\end{equation}


For expositional purposes, consider the following two-stage description of the compound environment.\footnote{The stage structure merely factorises the compound channel, which is a one-shot experiment. Preferences over the timing of uncertainty resolution are therefore abstracted from.}
  First, the action $a$ generates only
an interim record $i$ through $P$; no payoff is delivered at this stage.  Record $i$ selects the continuation
channel $K_i$, through which the same action generates the payoff $o$ and a further record $r$.  The agent does not
act again.  The interim record remains visible, so the final record is $(i,r)$.  Under $q$ and $P$, record $i$ is
\emph{reached} if $\sum_aq(a)P(i\mid a)>0$.  For every reached record, let $q_i(a):=q(a)P(i\mid a)/\sum_{a'}q(a')P(i\mid a')$ denote the posterior distribution
of the action.
\begin{description}[leftmargin=*,style=sameline]

\item[(A8) Recordwise sure-thing principle.]
For every $q\in\D(A)$, every binary first-stage record channel $P:A\to\D(\{1,2\})$ for which record $1$ is
reached, and every three continuation channels $K,L,M:A\to\D(O\times R)$,
\begin{equation}
 (q_1,K)\succeq(q_1,L)
 \quad\Longleftrightarrow\quad
 \bigl(q,P\triangleright(K,M)\bigr)
 \succeq
 \bigl(q,P\triangleright(L,M)\bigr).
 \label{eq:recordwise-sure-thing}
\end{equation}

Fix $q$ and a binary first-stage record channel $P$, and suppose record $1$ is reached, inducing posterior $q_1$.
The two compound environments share the continuation $M$ after record $2$ and differ only after record $1$, where
one continues with $K$ and the other with $L$.  The axiom says that the compounds are ranked exactly as the lottery
$q_1$ in channel $K$ and the same lottery in channel $L$ are ranked in their block environment.  The common
continuation on the other branch cannot affect the comparison.\footnote{This is a version of Savage's sure-thing principle
with the event given by an interim record \cite{Savage1954}.  Unlike the exogenous event in Savage's formulation,
reaching the record can be informative here: $P(1\mid a)$ may vary with the action.  This is why the conditional
comparison is evaluated at $q_1$.  If record $1$ is never reached, the axiom is silent, as on Savage's null
events.  When $P$ is an action-independent public coin, $q_1=q$, and
Axiom~\textup{(A8)} reduces to ordinary mixture independence with the coin retained in the record.}

\end{description}

% ===============================================================
\section{The representation theorem}\label{sec:theorem}
% ===============================================================

For $q\in\D(A)$ and $K:A\to\D(O\times R)$, let
\[
 p_{qK}(o,r):=\sum_a q(a)K(o,r\mid a)
\]
be the induced distribution of visible consequences, and define
\begin{equation}
 \I_{qK}(A;O,R)
 :=\sum_{a,o,r}q(a)K(o,r\mid a)
 \log\frac{K(o,r\mid a)}{p_{qK}(o,r)}.
 \label{eq:mi}
\end{equation}
Summands with $q(a)K(o,r\mid a)=0$ are set to zero, the logarithm has any fixed base greater than one, and expectations below are
taken under the joint distribution $q(a)K(o,r\mid a)$.

\begin{theorem}\label{thm:main}
For a family $\{\succeq_K\}_K$ as described above, the following are equivalent.
\begin{enumerate}[label=\textup{(\roman*)},leftmargin=*]
\item The family satisfies Axioms~\textup{(A1)--(A8)}.
\item There are a nonconstant function $u:O\to\mathbb R$ and a number $\lambda>0$ such that, for every nonempty
finite $A,R$, every $K:A\to\D(O\times R)$, and every $q,p\in\D(A)$,
\begin{equation}
 q\succeq_Kp
 \quad\Longleftrightarrow\quad
 V(q,K)\ge V(p,K),
 \qquad
 V(q,K):=\E_{qK}[u(O)]+\lambda\I_{qK}(A;O,R).
 \label{eq:representation}
\end{equation}
\end{enumerate}
Moreover, under these equivalent conditions, block-supported cross-channel comparisons are represented on the same
scale: for every finite block environment $\bigsqcup_{k\in I}K_k$, every two distinct blocks $i,j\in I$, and every
$q_i\in\D(A_i)$ and $q_j\in\D(A_j)$,
\[
 q_i^i\succeq_{\bigsqcup_{k\in I}K_k}q_j^j
 \quad\Longleftrightarrow\quad
 V(q_i,K_i)\ge V(q_j,K_j).
\]
\end{theorem}

We call a family satisfying \textup{(A1)--(A8)} \emph{trace-tempered}, and likewise the criterion that
represents it.  Its central restriction is that payoff and trace enter separately.  The agent may value good
outcomes and may value evidence about her action, but she cannot place an extra value on a good payoff being the
revealing consequence.  Thus, if under the same lottery two channels give the same expected utility from payoffs
and induce the same distribution of posteriors over actions, the agent is indifferent between them, regardless of
how payoffs and records are paired (Corollary~\ref{cor:matching}, Appendix~\ref{app:representation-consequences}).  The
axioms imply the separation.

\subsection{Proof sketch}

Appendix~\ref{app:main-proof} gives the details. For the moment, fix an action lottery $q$ assigning positive probability to every action.
Write $Y=(O,R)$ for the observed consequence and $q_y$ for the posterior
on the realised action given $Y=y$. Channels inducing the same joint
distribution of $(O,q_Y)$ can be obtained from one another by garbling
records while preserving payoffs. A6 therefore makes them indifferent.
Select a channel with an independent coin and record the coin's result.
This mixes the joint distributions. A8 gives mixture independence, so A1
and A2 yield a utility $W_q$ that is affine in these mixtures.

With one action, this gives expected utility over payoffs. A3 supplies
two strictly ranked payoffs to normalise its index $u$ by $u(o^-)=0$ and
$u(o^+)=1$. To compare different priors, draw their actions independently
and let each channel ignore the other draw. A7 makes this enlargement
irrelevant. At the common product lottery, affine uniqueness fixes every
$W_q$ in the same units. A5 ensures consistent comparisons across channels.

Next, change the continuation after a record of positive probability $m$.
A8 preserves the conditional ranking, so affine uniqueness makes the
ex-ante utility change a positive multiple of the conditional change.
Symmetry and comparisons through a record with a uniform posterior make payoff changes
under full revelation proportional to action probabilities, with a factor
that may depend on the prior. Comparing an upgrade directly and inside a
branch makes its multiplier $m$ times the ratio of the factors at the original and conditional
priors. Now form two equally likely cells, one carrying an arbitrary action
lottery and the other containing a single action. Report only the cell and
keep payoffs constant within it. A7 permits merging each cell into one action and reversing this by
redrawing within the cell. Compare changes from $o^-$ to $o^+$ in each cell
before and after merging. The conditional change in the cell with one action is worth one,
so its comparison equates the two experiments' factors. The other cell forces the factor for
its conditional action lottery to be one. Hence every branch multiplier is $m$.

To identify the information term, let $F(q)$ be the value of fully revealing the action at
payoff $o^-$, and $T_q(P)$ the value of a record channel $P$ at that payoff.
Write $m_y$ for its record probabilities and $q_y$ for its posteriors.
Appending full revelation on each reached branch adds $m_yF(q_y)$, so
\[
T_q(P)=F(q)-\sum_{y:m_y>0}m_yF(q_y).
\]
For a partition, A7 identifies the cell report with full revelation of
the cell. The identity gives the entropy grouping rule: full-revelation
value equals the value of learning the cell plus the average value of
learning the action within it. Comparisons with material lotteries give
continuity through A2. Together with symmetry, this allows Faddeev's theorem
to identify $F=\lambda H$, where $H$ is Shannon entropy.
A6 and A4 give $\lambda>0$.

Finally, run any channel $K$ by first recording its consequence $(o,r)$ and
then paying $o$. Replacing $o^-$ by $o$ after this record adds $u(o)$ times
the record's probability. Summing these changes gives
$V(q,K)=\E_{qK}[u(O)]+\lambda I_{qK}(A;O,R)$.
Deleting zero-probability actions removes the initial restriction on $q$.
The common normalisation gives the same $u$ and $\lambda$ across channels.

\subsection{Uniqueness}

A real-valued index $W$ on all finite lottery--channel pairs is a \emph{common-scale representation} if, for every channel
$K:A\to\D(O\times R)$ and every $q,p\in\D(A)$,
\[
 q\succeq_Kp
 \quad\Longleftrightarrow\quad
 W(q,K)\ge W(p,K),
\]
and, for every finite block environment $\bigsqcup_{k\in I}K_k$, every two distinct blocks $i,j\in I$, and all
$q_i\in\D(A_i)$ and $q_j\in\D(A_j)$,
\[
 q_i^i\succeq_{\bigsqcup_{k\in I}K_k}q_j^j
 \quad\Longleftrightarrow\quad
 W(q_i,K_i)\ge W(q_j,K_j).
\]

A common-scale representation $W$
is \emph{branch-additive} if, for every $k\ge1$, every finite $A$, every $q\in\D(A)$, every $P:A\to\D(\{1,\ldots,k\})$, and every
two ordered lists of continuation channels $(K_1,\ldots,K_k)$ and $(L_1,\ldots,L_k)$, with
$K_i,L_i:A\to\D(O\times R)$, writing
$m_i:=\sum_aq(a)P(i\mid a)$,
\begin{equation}
 W\bigl(q,P\triangleright(K_1,\ldots,K_k)\bigr)
 -W\bigl(q,P\triangleright(L_1,\ldots,L_k)\bigr)
 =\sum_{i:m_i>0}m_i
 \bigl[W(q_i,K_i)-W(q_i,L_i)\bigr].
 \label{eq:branch-additive-index}
\end{equation}
In words, a continuation change worth $x$ after record $i$ changes ex-ante value by $m_i x$, and changes on
different branches add. Axiom~\textup{(A8)} is ordinal: it says that a common continuation on the other branch cannot affect a conditional
comparison.  By Lemma~\ref{pt:lem:finite-branch-extension}\textup{(a)}, this is equivalent---given the other structural
axioms---to requiring that weak improvements at every reached branch give a weak compound improvement,
strictly if at least one reached branch improves strictly.  Branch additivity is the corresponding requirement
on numerical differences.

\begin{samepage}
\begin{proposition}[Uniqueness]\label{prop:uniqueness}
Suppose \textup{(A1)--(A8)} hold, and fix a representation $V$, with parameters $(u,\lambda)$, as in
Theorem~\ref{thm:main}.  For any real-valued index $W$ on all finite lottery--channel~pairs:
\begin{enumerate}[label=\textup{(\roman*)},leftmargin=*]
\item $W$ is a common-scale representation if and only if $W=\varphi\circ V$ for a single strictly increasing
$\varphi:[\underline u,\infty)\to\mathbb R$, where $\underline u:=\min_{o\in O}u(o)$.
\item A common-scale representation $W$ is branch-additive if and only if
$W=\alpha V+\beta$ for some $\alpha>0$ and $\beta\in\mathbb R$.
\end{enumerate}
\end{proposition}
\end{samepage}

The proof is in Appendix~\ref{app:representation-consequences}.

Every trace-tempered representation is a common-scale representation and, by the expected-utility and
mutual-information chain rules, branch-additive.  Part~\textup{(ii)} therefore implies that if $(u,\lambda)$ and
$(\widetilde u,\widetilde\lambda)$ are two such representations, then
\[
 \widetilde u=\alpha u+\beta,
 \qquad
 \widetilde\lambda=\alpha\lambda
\]
for some $\alpha>0$ and $\beta\in\mathbb R$.  Once two non-indifferent payoff outcomes are assigned numerical
utilities, $\lambda$ is unique for the chosen logarithm base.  On constant-payoff experiments alone, multiplying the trace index by a positive constant has no behavioural
content.  The material normalisation fixes the unit in which $\lambda$ measures the trade-off between payoff
and trace.

\subsection{Sign of the trace weight}

Axioms~\textup{(A4)}, \textup{(A6)}, and \textup{(A7)} orient the trace component.
Write $(\mathrm{A}4^-)$ for the reverse strict comparison in A4.  With the same quantifiers as in A6 and A7,
the reversed processing clauses are
\[
 (\mathrm{A}6^-):\quad(q,KT)\succeq(q,K),
 \qquad
 (\mathrm{A}7^-):\quad(qS,\widehat K)\succeq(q,K).
\]
Each uses the canonical block environment in the displayed order.  Write $(\mathrm{A}4^0)$,
$(\mathrm{A}6^0)$, and $(\mathrm{A}7^0)$ for indifference in the original comparison environments.

\begin{corollary}[Sign trichotomy]\label{cor:signduality}
Suppose Axioms~\textup{(A1)--(A3)}, \textup{(A5)}, and
\textup{(A8)} hold.  Then:
\begin{enumerate}[label=\textup{(\roman*)},leftmargin=*]
\item under \textup{(A4)}, \textup{(A6)}, and \textup{(A7)}, the family has the representation in
Theorem~\ref{thm:main}, with $\lambda>0$;
\item under $(\mathrm{A}4^-)$, $(\mathrm{A}6^-)$, and $(\mathrm{A}7^-)$, it has the same representation with
$\lambda<0$;
\item under $(\mathrm{A}4^0)$, $(\mathrm{A}6^0)$, and $(\mathrm{A}7^0)$, it is represented by expected utility,
equivalently by the same formula with $\lambda=0$.
\end{enumerate}
Conversely, each representation satisfies the corresponding orientation clauses.
In every case, the same parameters also represent finite block comparisons.
\end{corollary}

The proof is in Appendix~\ref{app:representation-consequences}.  The directions of the three orientation clauses cannot be chosen independently.  For example, under the
corollary's background axioms and any orientation of Axiom~\textup{(A7)}, Axiom~\textup{(A4)} is inconsistent
with $(\mathrm{A}6^-)$.  Restrict the two lotteries in Axiom~\textup{(A4)} to their supports and relabel the
actions.  The second channel is then a total garbling of the first.  Axiom~$(\mathrm{A}6^-)$ weakly prefers the
second, whereas Axiom~\textup{(A4)} strictly prefers the first.


\subsection{Independence of the axioms}

\begin{proposition}[Independence of the axioms]\label{prop:independence}
For each $i\in\{1,\ldots,8\}$ there is a family satisfying the other seven axioms and violating
Axiom~\textup{(A$i$)}.
\end{proposition}

The eight families are constructed in Appendix~\ref{app:axiom-independence}.  Three are of independent interest.  A conditional-entropy penalty satisfies every axiom except record data processing and values the removal of
extraneous record noise.  A Gini posterior-reduction criterion satisfies every axiom except action data processing.
The satiation criterion $\E[u(O)]+\min\{I(A;O,R),c\}$, for $c>0$, satisfies every axiom except the recordwise
sure-thing principle and ceases to value trace once $c$ units of information are secured.

% ===============================================================
\section{Choice implications}\label{sec:implications}
% ===============================================================

This section works out what an agent with the representation of Theorem~\ref{thm:main} does: which tests
falsify the model directly (\S\ref{sec:tracetest}), when she deliberately randomises (\S\ref{sec:randomisation}),
how $\lambda$ can be measured from choices (\S\ref{sec:measurement}), and how the criterion works in a
standard contracting model (\S\ref{sec:application}).
Throughout, after fixing the affine normalisation of $u$ and the logarithm base, $\lambda$ is measured in payoff units per unit of information.  Proofs for this
section, together with full statements of the supporting results that the text presents in summary form, are
collected in Appendix~\ref{app:choice-proofs}.

\subsection{The trace test}\label{sec:tracetest}

The following two-treatment test detects a preference for action--consequence dependence, distinguishing
it from a taste for randomisation alone or for the distribution of consequences alone.  It does not identify
the representation, whose quantitative restrictions are tested in \S\ref{sec:measurement}.

\begin{proposition}[The trace test]\label{prop:signature}
Four actions pay the same outcome $o_0$.  Two of them leave distinct marks: $a_1$ always produces record $r_1$,
and $a_2$ always produces the distinct record $r_2$.  The other two are indistinguishable: $a_3$ and $a_4$ each
produce $r_1$ or $r_2$ with equal probability.  Call this channel $K$, and let $q'$ be the fair coin over
$\{a_1,a_2\}$ and $q''$ the fair coin over $\{a_3,a_4\}$.
\begin{enumerate}[label=\textup{(\roman*)},leftmargin=*]
\item The two coins induce the same distribution of visible consequences, $p_{q'K}=p_{q''K}$; yet
\[
 V(q',K)-V(q'',K)=\lambda\log2>0,
 \qquad\text{so}\qquad
 q'\succ_Kq''.
\]
\item Let $K^\circ$ be the channel in which every action pays $o_0$ and produces $r_1$ or $r_2$ with equal
probability.  Then $q'\sim_{K^\circ}q''$.
\end{enumerate}
\end{proposition}


Part \textup{(i)} rules out any criterion that depends only on the distribution of visible consequences.
The two coins induce the same distribution over payoff--record pairs, so any such criterion is indifferent between them. However, a theory that also reads the lottery can
match it.  Existing theories of deliberate randomisation are of this kind: in the hedging model
of Cerreia-Vioglio, Dillenberger, Ortoleva, and Riella \cite{CerreiaVioglioDillenbergerOrtolevaRiella2019}, and
in the preference for randomisation documented by Agranov and Ortoleva \cite{AgranovOrtoleva2017}, the value of
mixing depends only on the induced distribution and the lottery.  Call an index $W$ \emph{channel-blind} if it
reads only these two objects: there is a function $F$, fixed across channels, such that
\[
 W(q,K)=F\bigl(p_{qK},q\bigr).
\]

Part \textup{(ii)}  rules these out.  Nothing a channel-blind index reads changes between the treatments: in
all four lottery--channel combinations the distribution of visible consequences is uniform on
$\{(o_0,r_1),(o_0,r_2)\}$, and the lotteries are the same two lotteries.  A channel-blind index therefore
separates the coins in both treatments or in neither, so none reproduces the choice pattern of
Proposition~\ref{prop:signature}.%
\footnote{If $W(q,K)=F(p_{qK},q)$ is channel-blind and $F(p,\cdot)$ is invariant under permutations of action labels---an entropy bonus, say---the first
treatment is enough: $p_{q'K}=p_{q''K}$ and $q''$ is a permutation of $q'$, so $W(q',K)=W(q'',K)$.  The second treatment is needed for label-sensitive indices.}

The first treatment also identifies the orientation of taste: it selects the revealing coin when $\lambda>0$,
the noisy coin when $\lambda<0$, and indifference when $\lambda=0$.  A strict preference between the coins in the second treatment rejects all three cases of Corollary~\ref{cor:signduality}, and with them the model.

\subsection{Deliberate randomisation}\label{sec:randomisation}

A qualitative preference for randomisation already follows from a subset of the axioms
(A1, A4--A7; Lemma~\ref{lem:qualitative-trace}).  Here we use the representation to show how the
trade-off between trace and material payoffs determines whether and how the decision maker randomises.


Consider a decision maker free to choose lotteries among two actions, and let $a$ pay weakly more.  Playing $a$ for sure maximises payoff and leaves no
trace; mixing in $b$ costs payoff but gives the record something to reveal.  Write $K_a:=K(\cdot\mid a)$ for the distribution of visible
consequences under action $a$, $\bar u(a):=\E_{K_a}[u(O)]$ for its expected utility, and
$\Delta:=\bar u(a)-\bar u(b)\ge0$ for the payoff gap.  Each unit of probability moved from $a$ to $b$ costs
$\Delta$.  What it buys is evidence.  A consequence is evidence about the act that produced it through a
likelihood ratio: how much more probable it is given the realised act than given only the mixture she plays.
With all probability on $a$, consequences are distributed exactly as $K_a$: every draw has ratio one and
carries no evidence.  Now move a marginal unit of probability to $b$.  The mixture's consequences are still,
to first order, distributed as $K_a$; but when $b$ realises, its consequence is drawn from $K_b$, so the draw
carries $\log\bigl(K_b(o,r)/K_a(o,r)\bigr)$.  Its expected value over $b$'s draws is, by definition, the
divergence $\KL(K_b\,\|\,K_a)$.\footnote{For distributions $P$ and $Q$,
$\KL(P\,\|\,Q):=\sum_{o,r}P(o,r)\log\bigl(P(o,r)/Q(o,r)\bigr)$, with zero-mass summands set to zero.  The
divergence is $+\infty$ if $P$ assigns positive probability to a consequence and $Q$ assigns zero.}  Because
information is concave in the lottery, whether to mix at all is settled at the margin (Proposition~\ref{prop:binary},
Appendix~\ref{app:randomisation-proofs}):
\[
 \delta_a\text{ is optimal}
 \quad\Longleftrightarrow\quad
 \Delta\ge\lambda\,\KL(K_b\,\|\,K_a).
\]
When the payoff gap is below $\lambda\KL(K_b\,\|\,K_a)$, the optimal lottery is unique and mixes the two
actions, with weakly more probability on $b$ as $\lambda$ rises---strictly when $\Delta>0$---and less as $\Delta$ rises.


With more than two actions, each action is compared with the distribution of visible consequences generated by
the lottery as a whole, $p_{qK}$. An action is credited with its expected utility and with how far its visible consequences depart from those generated on average.  At an optimum this total is the same for every action in use; an unused action cannot
offer more.

\begin{proposition}[Optimality condition]\label{prop:optimal-choice}
Let $\lambda>0$.  A lottery $q^*$ is optimal if and only if there is a constant $c$ such that
\[
 \bar u(a)+\lambda\KL\!\left(K_a\,\middle\|\,p_{q^*K}\right)=c
 \quad\text{for }a\in\supp(q^*),
 \qquad
 \le c\quad\text{for }a\notin\supp(q^*).
\]
\end{proposition}


The condition has an immediate implication: every visible consequence that can arise under some action
occurs with positive probability under any optimal lottery (Proposition~\ref{cor:support-pure},
Appendix~\ref{app:randomisation-proofs}).  Hence, if each action can produce some consequence that no other action can produce, every optimal lottery uses every action. At the other extreme, suppose every action produces the same distribution of payoffs and records. The consequence then reveals nothing about the action, and all actions have the same expected utility. Hence the agent is indifferent among all lotteries.\footnote{With $\lambda<0$, every payoff-maximising pure action is optimal.  More generally, $q$ is optimal in a fixed channel if and only if it is supported on payoff-maximising actions and $\I_{qK}(A;O,R)=0$.}


Now suppose that each action is fully identified by its visible consequences.  The next result shows that behaviour is then multinomial logit and the preference for trace becomes a preference for randomisation.

\begin{corollary}[Multinomial logit]\label{rem:logit}
If the row supports of $K$ are pairwise disjoint, then $\I_{qK}(A;O,R)=\Sh(q)$. Measuring information in natural logarithms, the unique optimum is
\[
 q^*(a)
 =\frac{\exp\bigl(\bar u(a)/\lambda\bigr)}
 {\sum_{a'}\exp\bigl(\bar u(a')/\lambda\bigr)},
\]
with $1/\lambda$ as the payoff sensitivity.\footnote{The same formula arises in rational inattention, where, in contrast, $q$ is fixed and $K$ is chosen \cite{MatejkaMcKay2015}.}
\end{corollary}

Here, Debreu's red-bus/blue-bus conclusion depends on what is recorded.  If blue is an exact copy of red, only their combined probability is pinned down. Any split of the red bus passengers between the two buses is optimal (Proposition~\ref{prop:optimal-marginal}, Appendix~\ref{app:randomisation-proofs}).  If colour is recorded, however, red and blue are distinct traces.  With equal expected utilities, the corollary assigns probability $1/3$ to each alternative.  The model is invariant to duplicate names, not to distinct recorded traces.


With overlapping supports, the odds between two actions can depend on the rest of the menu.  The optimality
condition sets each action's probability against the average consequence distribution, and an added action
shifts that average toward its own consequences, reducing the measured distinctiveness of the actions it
resembles more than of the others; the odds then move.  Pairwise-disjoint supports therefore imply independence
of irrelevant alternatives, while overlap permits, but does not require, its failure.


\subsection{Measuring the trace weight}\label{sec:measurement}

Under the stated normalisation, the coefficient $\lambda$ is the price of one unit of information in expected-utility units.  Any indifference
between two environments with different amounts of trace reveals it: writing $E_{qK}:=\E_{qK}[u(O)]$, if
$(q,K)\sim(p,L)$ and $\I_{qK}\ne\I_{pL}$, the representation rearranges to
\[
 \frac{E_{pL}-E_{qK}}{\I_{qK}-\I_{pL}}=\lambda,
\]
and $\lambda$ is the only number consistent with every such indifference
(Proposition~\ref{prop:single-price}, Appendix~\ref{app:measurement-proofs}).

Plot expected payoff utility against information: indifference curves are parallel lines of slope $-\lambda$.  The slope does
not depend on the prior, the payoff level, the action alphabet, or the amount of information already present.
This constancy is one observable restriction that Axiom~\textup{(A8)} adds to the rest of the system.  A taste for trace that satiates satisfies every axiom except~\textup{(A8)} (Proposition~\ref{prop:independence}) and has a slope that falls to zero beyond the satiation point.


The identity presupposes the utility function; the design in the next proposition does not, because its only
payoffs are the two normalised outcomes and a baseline that cancels.  At a constant payoff, revelation of the
realised action is worth $\lambda\Sh(q)$ more than silence, so compensating the silent side until the agent is
indifferent measures her willingness to pay for trace.

\begin{proposition}[Eliciting the trace weight]\label{prop:elicit}
Normalise $u(o^+)=1$ and $u(o^-)=0$.  Fix $q\in\D(A)$ with $\Sh(q)>0$, a constant payoff, and two channels, one
whose record reveals the realised action and one whose record is silent.  Dilute both by the same public side
branch, reached with probability $\varepsilon\in(0,1)$: on the revealing side the branch pays $o^-$ for sure; on
the silent side it pays $o^+$ with probability $\beta$ and $o^-$ otherwise.  If
$\lambda\Sh(q)\le\varepsilon/(1-\varepsilon)$, a unique $\beta\in[0,1]$ gives indifference,~and
\[
 \lambda=\frac{\varepsilon}{1-\varepsilon}\,\frac{\beta}{\Sh(q)};
\]
otherwise the revealing side is strictly preferred at every $\beta$.
\end{proposition}

The common baseline payoff's utility cancels, so it need not be known; the information gap $\Sh(q)$ is set by
the design, so no information quantity need be estimated.  The proof, the general compensation identity, and
the experimental controls are in Appendix~\ref{app:measurement-proofs} (Proposition~\ref{prop:wtp} and
Remark~\ref{rem:elicitation}).

The design also tests the model.  Run it twice with the same $\varepsilon$, at $q=(1/2,1/2)$ and at
$q=(1/4,3/4)$.  Two runs must return the same $\lambda$. Hence, the two indifference
reports satisfy
\[
 \frac{\beta_1}{\beta_2}
 =\frac{\Sh(1/2,1/2)}{\Sh(1/4,3/4)}
 =\frac{\log 2}{\frac14\log4+\frac34\log\frac43}
 \approx1.23.
\]
 Any other ratio is not consistent with
the model.  This number also separates Shannon entropy from the Gini-reduction rival
$1-\sum_aq(a)^2$, which instead commits to $4/3$.


% ===============================================================
\subsection{An application: moral hazard with a trace-tempered agent}\label{sec:application}
Take the simplest version of Holmstr\"om's moral-hazard problem \cite{Holmstrom1979}: a principal cannot observe
the agent's effort and offers a wage contingent on observable output.  The risk-neutral agent's effort is binary,
$A\in\{0,1\}$; high effort costs $c>0$ and low effort is costless.  Output $X\in\{0,1\}$ succeeds with
probability $p_a$ under effort $a$, where $0<p_0<p_1<1$, and the contract pays $W=w_0+\Delta wX$, with
$\Delta w\ge0$.  Write
$K_a$ for the Bernoulli row with parameter $p_a$, let $q:=\Pr(A=1)$, and attach the expected effort cost
$qc$ to the chosen lottery.\footnote{Writing cost at the level of the lottery prevents its realisation
from becoming a trace of effort.  The cost can nevertheless be kept inside the paper's finite consequence
domain.  For each candidate $q$, append to both action rows a payoff coordinate
$Z_q\in\{0,c\}$ with $\Pr(Z_q=c)=q$, and take $u(w,z)=w-z$.  Then $\E Z_q=qc$ and, because $Z_q$ is
independent of all production variables while $W$ is a function of $X$,
$\I(A;W,X,Z_q)=\I(A;X)$.  A plan $q$ is represented by the binary lottery $(1-q,q)$ in its augmented channel; any two
such plans are compared by placing their channels in separate blocks and applying the block-supported
common-scale clause of Theorem~\ref{thm:main}.}

For a given contract $(w_0, \Delta w)$,  the agent chooses $q\in[0,1]$ to maximise
\[
 V(q)=w_0+\bigl[p_0+q(p_1-p_0)\bigr]\Delta w-qc
      +\lambda\,\I_{qK}(A;X).
\]
Pure high effort is optimal if and only if
\[
 \Delta w\ge \frac{c+\lambda\,\KL(K_0\,\|\,K_1)}{p_1-p_0}
\]
(Appendix~\ref{app:moral-hazard-proofs}). Thus the performance bonus must cover both effort cost and the trace forgone
at a pure action.  At a flat wage, the unique optimum is interior if and only if
\[
 c<\lambda\,\KL(K_1\,\|\,K_0),
\]
so trace alone can sustain some effort without an incentive bonus.

Now consider the team setting from \cite{Holmstrom1982}. Suppose $n$ agents choose efforts and the output is $X\in\{0,1\}$  with success probability $\varepsilon < g(a_1,\ldots,a_n) < 1-\varepsilon$ for some $\varepsilon>0$. Fix the
strategies of agent $i$'s teammates.  If $i$ works, the signal is positive with probability
$\bar g_1:=\E[g(1,A_{-i})]$; if she shirks, with probability $\bar g_0:=\E[g(0,A_{-i})]$.

Assume positive marginal productivity, $\bar g_1>\bar g_0$. The agent's problem is then the one described above with $p_a$ replaced by $\bar g_a$, and both conditions apply as written.  The minimum bonus that sustains pure high effort is now
\[
 \Delta w^*=\frac{c+\lambda\,\KL(\bar K_0\,\|\,\bar K_1)}{\bar g_1-\bar g_0},
\]
where $\bar K_a$ is the Bernoulli row with parameter $\bar g_a$.

In technologies in which one member's effect vanishes as the team grows---for example, when success is a uniformly Lipschitz function of the fraction working---success is almost as likely whether she works or shirks, and a draw of $X$
says almost nothing about her effort.
The gap $\bar g_1-\bar g_0$ then closes.  The effort share of the bonus grows as it does in the standard
model; the divergence is second order in the gap, so the trace share vanishes: trace becomes second order
relative to pay.  The exception is an essential task, whose effect does not shrink with the team.

Suppose a signal $S$ of her work is visible---posted within the team, say---but not verifiable, so wages
remain contingent on $X$ alone.  Recording it beside $X$, at $q$ held fixed, raises her value by
$\lambda\,\I_q(A_i;S\mid X)$ (Proposition~\ref{prop:garbling-price}, Appendix~\ref{app:record-value-proofs}).
If the base wage can be adjusted freely, the principal can reduce it by the increase in the agent's maximised utility while preserving participation. This increase is at least $\lambda\,\I_{q^*}(A_i;S\mid X)$ at any original optimum $q^*$, but the added signal may change effort. Non-contractible measurement can therefore substitute for base pay; Appendix~\ref{app:moral-hazard-proofs} gives the exact compensation formula. With $\lambda=0$, this direct value is zero.


% ===============================================================
\section{Related literature}\label{sec:literature}
% ===============================================================

Grant, Kajii, and Polak \cite{GrantKajiiPolak1998} characterise an intrinsic preference for information
independent of its use in planning; Caplin and Leahy \cite{CaplinLeahy2001} model anticipatory feelings generated
by beliefs, while Ely, Frankel, and Kamenica \cite{ElyFrankelKamenica2015} study preferences over paths of beliefs
and the resulting value of suspense and surprise.  These papers show that information need not be instrumental,
but the relevant beliefs concern uncertainty faced by the agent, whereas the posterior here concerns her
realised action.  The axioms here also restrict its value to expected entropy reduction and fix its rate of
substitution with material utility.

Cerreia-Vioglio, Dillenberger, Ortoleva, and Riella
\cite{CerreiaVioglioDillenbergerOrtolevaRiella2019} study deliberate randomisation and characterise a model in
which mixing hedges uncertainty about how to evaluate lotteries; Agranov and Ortoleva
\cite{AgranovOrtoleva2017} document stochastic choice under announced repetition and evidence that much of it is
deliberate.  Trace-tempered choice can also make a nondegenerate lottery optimal. The test in Proposition~\ref{prop:signature} tells the two models apart.

Corrao, Fudenberg, Mudekereza, and Levine \cite{CorraoFudenbergMudekerezaLevine2026} represent
concave and smooth utilities over lotteries as expected utility plus suspense, the expected error of an
adversary's best forecast of the outcome.  Within a fixed channel the present criterion has this form: the
adversary forecasts the visible consequence, and the suspense is $\lambda\,\I_{qK}(A;O,R)$. Under logarithmic loss, this is the excess prediction error from not knowing the action, after subtracting the error that remains when the action is known (Appendix~\ref{app:interpretation-proofs}). Their error function is a primitive. Here the axioms determine this information term, in every channel, with the same $\lambda$.

In signaling models, actions convey information about hidden characteristics.  Andreoni and Bernheim
\cite{AndreoniBernheim2009} let chance override the dictator's choice in a dictator game with known probability, and fair behaviour
falls: the observer can no longer tell a selfish choice from chance.  In B\'enabou and Tirole
\cite{BenabouTirole2006}, the agent values the audience's belief that her motives are good; a paid good deed
proves less, so rewards can crowd out good conduct.  In their earlier paper \cite{BenabouTirole2002} the
audience is internal: a future self whose confidence sustains effort.  Bernheim \cite{Bernheim1994} shows how
status concerns can induce agents with heterogeneous preferences to choose the same action, making behaviour less
informative about predispositions---an instrumental counterpart to the deniability case in
Corollary~\ref{cor:signduality}.  In all these models the audience infers a hidden characteristic from the
action, and the response depends on what is inferred: a better type earns more esteem.  Here no response enters
utility, and inference runs the other way, from consequence to action. The audience interpretation developed here pays for the amount learned, whatever its content (\S\ref{sec:interpretation}).

The proof also relates to axiomatic characterisations of entropy.  Shannon \cite{Shannon1948} characterised
entropy from continuity, monotonicity on uniform distributions, and a grouping recursion; Faddeev
\cite{Faddeev1956} reduced these to symmetry, continuity of the binary entropy, and the recursion itself.  The
proof uses this binary-continuity form of Faddeev's theorem; Baez, Fritz, and Leinster
\cite{BaezFritzLeinster2011} give a related strong-additivity formulation.
In those characterisations, a numerical information functional is primitive and structural conditions are
imposed directly on it. Here the information chain rule and Faddeev's recursion are derived from axioms on an ordinal family of weak orders over action lotteries. The theorem also calibrates the information term against material utility.

\looseness=-1 Blackwell \cite{Blackwell1953} compares experiments by their value in all downstream decision problems; Gilboa
and Lehrer \cite{GilboaLehrer1991} characterise which functions on partitions can arise as values of information, and Azrieli and Lehrer \cite{AzrieliLehrer2008} extend the analysis to stochastic
information structures.  Frankel and Kamenica \cite{FrankelKamenica2019} characterise the measures of
information and uncertainty that arise as the value of news.  Cabrales, Gossner, and
Serrano \cite{CabralesGossnerSerrano2013} obtain entropy reduction from the decisions of ruin-averse investors.
For Blackwell's decision-maker, information is an input: she can use it or ignore it, so more
information cannot hurt.  Here it is an output: the record is evidence about the agent, and no
decision follows it.
Uniform preference between two record technologies identifies the more-capable
order of K\"orner and Marton \cite{KornerMarton1977}, not Blackwell garbling
(Proposition~\ref{prop:blackwell}, Appendix~\ref{app:record-value-proofs}).

Mutual information is also the cost of attention in rational inattention \cite{Sims2003,MatejkaMcKay2015}. The difference is in what is fixed and what is chosen. There the distribution of the state is fixed and the agent chooses how her action depends on it. The cost is the information the action carries about the state.  Here the channel from actions to consequences is fixed and the agent chooses the distribution of her action. The benefit is the information the consequence carries about
the action. The data that provide identification differ as
well.  Rational inattention is identified from state-dependent stochastic choice; Denti \cite{Denti2022} gives
testable conditions under which such data are rationalised by posterior-separable costs.  Here the primitive
is a ranking of action lotteries in a fixed channel, and the trace test compares two lotteries that induce the
same distribution of consequences.  When the rows of the channel have disjoint supports, both models produce
choice probabilities of the multinomial-logit form, so choice frequencies alone do not separate them.
Rankings of lotteries that are not chosen and comparisons across channels do.  Caplin, Dean, and Leahy
\cite{CaplinDeanLeahy2022} show that, among uniformly posterior-separable costs, Invariance under Compression
selects Shannon entropy: merging states that pay the same in every action does not change behaviour.  The
corresponding condition here is indifference to undetectable distinctions, which follows from action data
processing: merging two actions with identical rows, and conditionally redrawing within the merged cell,
are opposite processors under Axiom~\textup{(A7)}, so neither changes the experiment's value.  The Gini criterion in Appendix~\ref{app:axiom-independence} satisfies every
axiom except action data processing and fails it at such a merge.

Pomatto, Strack, and Tamuz \cite{PomattoStrackTamuz2023} ask what an experiment costs.  On a cardinal,
prior-free cost function they impose continuity, monotonicity in Blackwell's order, additivity over
independent experiments, and constant marginal cost: an experiment run with probability $\alpha$ costs
$\alpha$ times as much.  Then the only answer is a weighted sum of Kullback--Leibler divergences
between the signal distributions in different states.
Their cost functional and the trace value studied here do not nest.
Full revelation has trace value $\Sh(q)$, finite and prior-dependent, while every nonzero cost in their class
makes it infinitely costly.  The information-cost literature prices
learning about the world; the present paper prices being learned about.


% ===============================================================
\section{Interpretation and scope}\label{sec:interpretation}
% ===============================================================

\paragraph{Evidence, not authorship or control.}
The trace term prices evidence of agency: not authorship, and not control in its everyday sense.  A pure action
fixes the act in advance, although its consequence may remain stochastic; it has
zero trace value even when the consequence names the act perfectly, because an act settled in advance leaves no
uncertainty for the consequence to resolve.  A coin the agent flips herself can do better.  The model does not
say that surrendering a choice to chance creates more agency; it says that an unsettled act can leave more
evidence than a predetermined one.  The primitive ranks lotteries she implements herself; a lottery assigned to her
from outside is not in the domain.  This is the Underground Man's problem: the proof he wants cannot come from an act
fixed in advance.  A reading of control or authorship survives if what is valued is its display rather than its
possession; but the display of which act was taken is exactly what the trace term measures, so the surviving reading
is the trace reading.

\paragraph{Relation to empowerment.}
The channel maps actions into distributions over consequences, and the information term can be read from either
end.  Read backward, from consequence to action, it measures trace: how much the visible consequence reveals
about the realised action.  Read forward, from action to consequence, it measures how strongly the action
shapes the distribution of visible consequences.  Mutual information is symmetric, so the two readings share
one index.  No choice among lotteries separates them.

The forward reading has an exact counterpart.  Suppose $\E[u(O)\mid a]=\bar u$ for every action.  Since
$\lambda>0$, maximising $V$ over $q$ gives
\[
 \max_{q\in\Delta(A)}V(q,K)=\bar u+\lambda C(K),
 \qquad
 C(K):=\max_{q\in\Delta(A)}I_{qK}(A;O,R),
\]
the Shannon capacity of the channel.  When the visible consequence is read as the agent's sensor state, $C(K)$
is the one-step empowerment index of Klyubin, Polani, and Nehaniv \cite{KlyubinPolaniNehaniv2005}; see also
\cite{SalgeGlackinPolani2014}.  Under this supposition, the trace-tempered agent solves exactly the problem
that defines one-step empowerment: her optimal lotteries are the capacity-achieving action distributions, and
the size of $\lambda$ plays no role in the ranking---the price is visible only against payoffs.  Empowerment
summarises the channel by the optimised value $C(K)$; the representation recovers the ranking beneath it,
$q\mapsto I_{qK}(A;O,R)$, from preferences and prices it against material utility.

\paragraph{What the axioms preclude.}
The main exclusion is selective legibility.  Mutual information responds to an act's
probability and to the distinctiveness of its consequences, not to what the act is.  An agent who wants her generous
acts recognised and her mean acts hidden may therefore prefer a channel to its garbling in one case and the garbling to its source
in another.  Such sign-reversing preferences violate the global data-processing requirement in Axiom~\textup{(A6)} and fall
outside all three cases of the sign trichotomy.  Such preferences are not irrational; they are tastes for
something other than trace.  Wanting particular acts known points to an audience, outside the agent or within
her, for whom what is revealed matters.  The taste axiomatised here does not depend on one: it prices the amount
the consequence reveals.  A model of selective legibility would have to supply the audience, making
the value of evidence depend on its content and not only on the amount of uncertainty it removes.

\paragraph{The payoff--record distinction.}



Although the information term treats $O$ and $R$ jointly, the axioms give them different roles.  In particular, Axiom~\textup{(A6)} permits the record to be garbled while the payoff is held fixed. This separates a change in evidence from a change in material
consequences.

 Axiom~\textup{(A8)} permits records to arrive in stages while the payoff is delivered once.  Its recordwise sure-thing restriction makes branchwise comparisons commensurable and supplies the cardinal structure used to place payoff utility and information on a common scale.  Thus either $O$ or $R$ may reveal the action, but the axioms imply that only
$O$ enters material utility.  A record
that is itself pleasant or shameful belongs in $O$.  Provided the complete visible information is preserved, this reclassification leaves mutual information unchanged.




\paragraph{One shot, with no audience.}
The model describes one episode.  It specifies neither an audience nor any response to the beliefs induced by
the record; the posterior is a statistic of the implemented lottery and the channel, not anyone's belief.  The
stages in Axiom~\textup{(A8)} factor a single one-shot interaction: no action follows the record, and no relationship
extends beyond the episode.  Now take the instrumental reading of the introduction literally: an observer
watches the visible consequence and rewards the agent at rate $\lambda$ for the fall in the entropy of its beliefs about her action. Measure the reward in utility units, or in money under quasilinearity. Its expectation is $\lambda\,\I_{qK}(A;O,R)$ when the observer's prior is the implemented lottery $q$; only this prior gives the equality for every channel (Proposition~\ref{prop:audience}, Appendix~\ref{app:interpretation-proofs}). To reproduce rankings over plans, its prior must track each implemented lottery; correctness at one equilibrium lottery is insufficient for deviations against a fixed prior. The audience cares about how much is revealed, not which act it was, so selective legibility remains excluded.  The audience
is not a primitive: the representation predicts no change when an audience that pays nothing is added while the action
lottery and channel are held fixed.  If choices change, the model is missing the value of being observed by
others.

\clearpage
\appendix

\part*{Proofs}
\addcontentsline{toc}{part}{Proofs}

\section{Proof of the representation theorem}
\label{app:main-proof}

We prove Theorem~\ref{thm:main} by first fixing a common material utility
scale.  This scale identifies the numerical weight of a reached continuation
as its probability, before any information formula is derived.  Appending
full revelation then gives an entropy-reduction identity; action processing
and continuity identify its entropy as Shannon's.  Deterministic payoff
separation and sequentialisation complete the representation.
The preliminary lemmas state their own hypotheses.  Thereafter we assume
Axioms~\textup{(A1)--(A8)} until the converse proof at the end of the appendix.

\subsection{Shared comparison and affine tools}

Throughout the proof, comparisons between experiments are the canonical
two-block comparisons of the main text.  We first record how these comparisons
behave under changes of alphabet and under continuation replacement.

\begin{lemma}[Blocks, supports, and dummies]\label{lf:lem:bookkeeping}
Under Axioms \textup{(A1)} and \textup{(A5)--(A7)}:
\begin{enumerate}[label=\textup{(\alph*)},leftmargin=*]
\item The canonical comparison of experiments is complete and transitive,
and is unchanged by invertible action or record relabellings and by
replacement of either experiment with an indifferent one.
\item If $S=\supp q$, $q_S$ is the restriction of $q$ to $S$, and $K_S$
is the restriction of $K$ to those action rows, then
$(q_S,K_S)\sim(q,K)$.
\item If $s\in\D(B)$ and
$K^\uparrow(o,r\mid a,b)=K(o,r\mid a)$, then
\begin{equation}
(q\otimes s,K^\uparrow)\sim(q,K).
\label{lf:eq:dummy-ordinal}
\end{equation}
\end{enumerate}
\end{lemma}

\begin{proof}
Place a finite collection of experiments in one block environment.
Axiom~\textup{(A5)} removes unused blocks and Axiom~\textup{(A1)} gives its
weak order.  Invertible relabellings are neutral by two-way applications of
Axioms~\textup{(A6)} and~\textup{(A7)}; in particular, swapping the two
block tags identifies the two orientations of a comparison.  Three blocks
give transitivity and four give replacement of indifferent experiments.

For support restriction, use the inclusion $S\hookrightarrow A$ in one
direction and a retraction $A\to S$ in the other.  For dummy extension,
project $(a,b)$ to $a$ in one direction and, in the other, draw $b$ independently
from $s$ after observing $a$.  Each pair of processors satisfies
\eqref{eq:action-processing} in both directions; off-support rows admit
arbitrary completions.  Axiom~\textup{(A7)} proves the asserted indifferences.
\end{proof}

\begin{lemma}[Finite-branch extension]\label{pt:lem:finite-branch-extension}
Suppose Axioms~\textup{(A1)} and \textup{(A5)--(A7)} hold.
\begin{enumerate}[label=\textup{(\alph*)},leftmargin=*]
\item Axiom~\textup{(A8)} is equivalent to its following finite-branch form.
For a first-stage channel $P:A\to\D(\{1,\ldots,k\})$, prior $q$, and
continuations $K_i,L_i:A\to\D(O\times R_i)$, put
\[
m_i=\sum_a q(a)P(i\mid a),\qquad
q_i(a)=\frac{q(a)P(i\mid a)}{m_i}\quad(m_i>0).
\]
If $(q_i,K_i)\succeq(q_i,L_i)$ at every reached branch, then
\begin{equation}
\bigl(q,P\triangleright(K_i)_i\bigr)
\succeq\bigl(q,P\triangleright(L_i)_i\bigr),
\label{pt:eq:finite-branch-extension}
\end{equation}
strictly if a reached branch comparison is strict.  The two continuations at
each branch may have a common branch-dependent record alphabet.
\item Under Axiom~\textup{(A8)}, if the lists agree except at a reached
branch $i^*$, then
\[
(q_{i^*},K_{i^*})\succeq(q_{i^*},L_{i^*})
\quad\Longleftrightarrow\quad
\bigl(q,P\triangleright(K_i)_i\bigr)
\succeq\bigl(q,P\triangleright(L_i)_i\bigr),
\]
and likewise for strict preference.
\end{enumerate}
\end{lemma}

\begin{proof}
Embed the finitely many record alphabets as tagged subsets of one finite
alphabet.  The tags can be added and deleted by payoff-preserving record
processors.  If $m_i=0$, then $P(i\mid a)=0$ on $\supp q$, so changing that
continuation changes only irrelevant action rows.  If just one branch is
reached, its first-stage label is constant and can be deleted and restored.
Thus consider at least two reached branches.

Replace their continuations one at a time.  At the step replacing branch
$i$, use the binary first stage
$B(1\mid a)=P(i\mid a)$ and $B(2\mid a)=1-P(i\mid a)$.
On branch $1$, keep $K_i$ or $L_i$ and append the tag $i$.
On branch $2$, keep the common continuation
\begin{equation}
M(o,(j,r)\mid a)
=\frac{P(j\mid a)}{B(2\mid a)}H_j(o,r\mid a),\qquad j\ne i,
\label{pt:eq:complement-continuation}
\end{equation}
whenever the denominator is positive, with arbitrary completion otherwise;
here $H_j$ is the unchanged continuation at $j$.  The posterior at binary
branch $1$ is exactly $q_i$.  Deleting the redundant binary tag identifies
each binary compound with the corresponding original compound: on the active
record support, $(1,(i,r))\mapsto(i,r)$ and
$(2,(j,r))\mapsto(j,r)$ are invertible record maps.
Axiom~\textup{(A8)} therefore ranks this replacement weakly, and strictly
when its branch comparison is strict.  Transitivity in a common finite-block
environment proves part~\textup{(a)} in the forward direction.

Conversely, apply the finite-branch property to the binary lists in
Axiom~\textup{(A8)}, whose other continuation is common.  If a weak aggregate
comparison held while its branch comparison failed, completeness would give
a strict reverse branch comparison, and the strict finite-branch property
would give a strict reverse aggregate comparison.  This proves the reverse
implication in Axiom~\textup{(A8)}.  The same argument proves
part~\textup{(b)}; applying its weak equivalence in both directions gives
the strict equivalence.
\end{proof}

Below, finite-branch uses of Axiom~\textup{(A8)} refer to this lemma.
In every compound the action is drawn only once; $q_i$ is its posterior
after the first-stage record.

\subsection{Marked terminal laws and affine utility}

For experiments at a common prior $q$, write $E\succeq_q G$ for
$(q,E)\succeq(q,G)$.  At a full-support prior $q\in\D(A)$, define
\[
p_K(o,r)=\sum_a q(a)K(o,r\mid a),\qquad
\pi^K_{or}(a)=\frac{q(a)K(o,r\mid a)}{p_K(o,r)}
\quad\text{when }p_K(o,r)>0.
\]
The \emph{marked terminal law} is the joint distribution of the payoff and
the posterior on the realised action:
\begin{equation}
\eta_K=\sum_{(o,r):p_K(o,r)>0}
p_K(o,r)\,\delta_{(o,\pi^K_{or})}.
\label{lf:eq:marked-law}
\end{equation}
Its posterior marginal has barycentre $q$.  We sometimes write $\eta_{q,K}$
to display the prior explicitly.

\begin{lemma}[Terminal laws and affine fibres]\label{lf:lem:terminal-affine}
Under Axioms \textup{(A1)}, \textup{(A2)}, and \textup{(A5)--(A8)}:
\begin{enumerate}[label=\textup{(\alph*)},leftmargin=*]
\item At a full-support prior, experiments with the same marked terminal
law are indifferent.
\item At every full-support $q$, the order on finite marked terminal laws
has an affine real-valued representation $W_q$.
\end{enumerate}
\end{lemma}

\begin{proof}
For part~\textup{(a)}, let
$M_K(o,\pi)=\sum_{r:\,p_K(o,r)>0,\,\pi^K_{or}=\pi}p_K(o,r)$.
Equality of marked laws gives $M_K(o,\pi)=M_L(o,\pi)=M(o,\pi)$.
For a reached output $(o,r)$ of $K$, define
\begin{equation}
T(s\mid o,r)=
\1_{\{p_L(o,s)>0,\,\pi^L_{os}=\pi^K_{or}\}}
\frac{p_L(o,s)}{M(o,\pi^K_{or})},
\label{lf:eq:matching-processor}
\end{equation}
and complete unreached rows arbitrarily.  This processor preserves the
payoff.  For a reached $(o,s)$ of $L$, with posterior $\pi$,
Bayes' formula gives
\[
\sum_r K(o,r\mid a)T(s\mid o,r)
=\frac{p_L(o,s)}{M(o,\pi)}\frac{\pi(a)}{q(a)}M(o,\pi)
=L(o,s\mid a).
\]
For an unreached target output both sides are zero, since $q$ has full
support.  Thus $K$ garbles to $L$ by record processing, and the symmetric
construction garbles $L$ to $K$.  Axiom~\textup{(A6)} gives indifference.

For part~\textup{(b)}, quotient experiments by equality of marked laws.
The public mixture \mbox{$tE\oplus(1-t)G$} selects an experiment using an independent
coin and retains that coin in the record.  Hence
\begin{equation}
\eta_{tE\oplus(1-t)G}=t\eta_E+(1-t)\eta_G.
\label{lf:eq:public-mixture}
\end{equation}
The action-independent first-stage case of Axiom~\textup{(A8)} gives
\begin{equation}
E\succeq_q L
\quad\Longleftrightarrow\quad
tE\oplus(1-t)G\succeq_q tL\oplus(1-t)G,
\qquad 0<t<1.
\label{lf:eq:mixture-independence}
\end{equation}
Different record alphabets are placed in a tagged disjoint union;
part~\textup{(a)} makes the resulting mixture well defined on the quotient.
A mixture segment and a fixed comparator can be realised on one fixed
finite record alphabet, with channel entries continuous in $t\in[0,1]$.
Axiom~\textup{(A2)} makes both comparison sets on that segment closed.
If the comparator lies between the endpoints, completeness makes those sets
cover the segment, and connectedness makes them intersect.  This is the
Archimedean condition.  The mixture-space theorem
\cite{HersteinMilnor1953} supplies the affine representation.
\end{proof}

\subsection{Relevance and a common material scale}

Write $K_o$ for the singleton-action experiment delivering $o$ for sure,
$K_o^{\mathrm{id}}$ for full revelation at payoff $o$, and $K_o^0$ for
the same payoff with a constant record.

\begin{lemma}[Fixed-channel and environment relevance]\label{lem:relevance-bridge}
Under Axioms~\textup{(A1)} and \textup{(A5)--(A8)}:
\begin{enumerate}[label=\textup{(\alph*)},leftmargin=*]
\item Axiom~\textup{(A3)} is equivalent to the existence of $o^+,o^-\in O$
such that $(\delta_*,K_{o^+})\succ(\delta_*,K_{o^-})$.
\item Axiom~\textup{(A4)} is equivalent to the existence of $o_*\in O$
such that $(q,K_{o_*}^{\mathrm{id}})\succ(q,K_{o_*}^0)$ for every
full-support $q$ on every finite action set with at least two elements.
\end{enumerate}
\end{lemma}

\begin{proof}
Two-way action processing identifies each pure lottery in Axiom~\textup{(A3)}
with its singleton experiment, proving part~\textup{(a)}.  For
part~\textup{(b)}, restrict the A4 benchmark lotteries to their supports.
Two-way action and record processing identify their comparison with
\begin{equation}
((\tfrac12,\tfrac12),K_{o_*}^{\mathrm{id}})
\succ((\tfrac12,\tfrac12),K_{o_*}^0).
\label{eq:bridge-fair-binary}
\end{equation}
At any full-support binary prior $p$, choose
$0<\varepsilon<\min\{1,2p_1,2p_2\}$ and a binary record channel with
$P(y\mid i)=\varepsilon/(2p_i)$.  Record $y$ has mass $\varepsilon$ and
fair posterior.  Revealing only after $y$ is strictly better than remaining
silent after both records, by \eqref{eq:bridge-fair-binary} and
Lemma~\ref{pt:lem:finite-branch-extension}.  Full revelation garbles to the
first compound, and the second garbles to silence, so full revelation is
strictly better than silence at $p$.
For general full-support $q$, reveal a partition into two nonempty cells.
Projection onto the cells and conditional redrawing within each cell are
opposite A7 processors identifying this experiment with binary full
revelation; silence is likewise identified with binary silence.  Full
revelation on the original alphabet dominates the cell report by A6.
This proves the desired strict comparison.  Conversely, apply the condition
to the fair binary prior and reverse the benchmark identifications.
\end{proof}

For a payoff lottery $\ell\in\D(O)$, write $\ell$ also for the experiment
that draws this payoff independently of the action and has a constant record.

\begin{lemma}[Material utility and a common affine scale]\label{lf:lem:material-scale}
Under Axioms \textup{(A1)--(A3)} and \mbox{\textup{(A5)--(A8)}}, there are
$o^+,o^-\in O$ and a nonconstant $u:O\to\R$, with $u(o^+)=1$ and
$u(o^-)=0$, such that the affine representatives can be normalised to satisfy
\begin{equation}
W_q(\ell)=\sum_o\ell(o)u(o)
\label{lf:eq:payoff-normalisation}
\end{equation}
for every full-support $q$ and payoff lottery $\ell$.
This normalisation gives exact dummy invariance
$W_{q\otimes s}(E^\uparrow)=W_q(E)$ at full-support priors and represents
every full-support comparison on a common scale:
\begin{equation}
(q,E)\succeq(q',G)
\quad\Longleftrightarrow\quad W_q(E)\ge W_{q'}(G).
\label{lf:eq:common-scale}
\end{equation}
\end{lemma}

\begin{proof}
On a singleton action set, the marked law is just the payoff lottery.
Lemma~\ref{lf:lem:terminal-affine} gives an affine representation there,
and Lemma~\ref{lem:relevance-bridge}\textup{(a)} supplies two strictly
ranked deterministic payoffs.  Normalise them to one and zero and let $u(o)$
be the value of the degenerate lottery at $o$.  Affinity gives the expectation
of $u$ on every payoff lottery.

An action processor can ignore its input and draw any target prior.
For an action-independent payoff lottery, this satisfies the A7 completion
equation in both directions.  The induced order on payoff lotteries is
therefore identical at every prior and on every action alphabet.
Its embedding into each full-support marked-law fibre is affine and
order-reflecting.  Positive-affine uniqueness, with the same two anchors,
gives \eqref{lf:eq:payoff-normalisation}.

Dummy lifting is also affine on marked laws: it sends $(o,\pi)$ to
$(o,\pi\otimes s)$.  By Lemma~\ref{lf:lem:bookkeeping}, its pullback
represents exactly the original fibre order.  Thus affine uniqueness gives
$W_{q\otimes s}(E^\uparrow)=cW_q(E)+d$, with $c>0$.
Both payoff anchors have the same values before and after lifting, so
$c=1$ and $d=0$.  The same argument applies to a lift that ignores the first
coordinate.  Finally lift $E$ at $q$ and $G$ at $q'$ to the same prior
$q\otimes q'$, with each experiment ignoring the other coordinate.
Replacement of indifferent experiments identifies their canonical comparison
with this fixed-prior comparison; dummy invariance then gives
\eqref{lf:eq:common-scale}.
\end{proof}

\paragraph{Boundary convention and exact invariances.}
For an arbitrary prior $q$ on $A$, henceforth define
\begin{equation}
W_q(E):=W_{q_S}(E_S),\qquad S=\supp q.
\label{lf:eq:support-normalisation}
\end{equation}
The representatives on the right have the material normalisation just fixed.
Support indifference in Lemma~\ref{lf:lem:bookkeeping} extends
\eqref{lf:eq:common-scale} to \emph{all} priors, including singletons.
Consequently every indifference established by two-way action or record
processing is an equality of these numerical values.  In particular,
relabelling actions and their priors, relabelling records, adjoining an
independent dummy, and deleting or reinstating zero-probability actions
preserve $W$ exactly.  When supports are embedded in larger alphabets,
posteriors are padded by zeros; this changes neither the order nor its
material-normalised value.  We retain the notation $\eta_{q,E}$ for the
marked law on the stated alphabet and omit zero-mass outputs throughout.

\subsection{Material calibration of continuation weights}
\label{mf:sec:calibration}

We now use the material unit to identify the weight of a continuation before
identifying the value of information.  Throughout this subsection assume
Axioms~\textup{(A1)--(A8)}, and use the common normalisation of
Lemma~\ref{lf:lem:material-scale}.  In particular, the uninformative payoffs
$o^-$ and $o^+$ have values zero and one at every prior, including a
singleton prior.

We first record a consequence of affinity that allows us to compare payoff
changes implemented in different ways.  Differences of marked laws below are
finite signed measures on $O\times\D(A)$.

\begin{lemma}[Equal changes in marked laws]\label{mf:lem:scaled-differences}
Fix a full-support prior $q$.  If four experiments at $q$ satisfy
\[
 s(\eta_E-\eta_G)=t(\eta_L-\eta_M),\qquad s,t>0,
\]
then
\[
 s[W_q(E)-W_q(G)]=t[W_q(L)-W_q(M)].
\]
\end{lemma}

\begin{proof}
Put $\theta=s/(s+t)$.  The public mixtures
$\theta E\oplus(1-\theta)M$ and
$\theta G\oplus(1-\theta)L$ have the same marked law.  Their values are
equal by Lemma~\ref{lf:lem:terminal-affine}; affinity and multiplication by
$s+t$ give the result.
\end{proof}

Let $q\in\D(A)$ have full support and let $P:A\to\D(Y)$ be a first-stage
record channel.  Write
\[
 m_y=\sum_a q(a)P(y\mid a),\qquad
 \widehat r_y(a)=\frac{q(a)P(y\mid a)}{m_y}\quad(m_y>0).
\]
For a reached $y$, put $S_y=\supp\widehat r_y$ and let $r_y$ be the
full-support restriction of $\widehat r_y$ to $S_y$.  Let
$\iota_y:\D(S_y)\to\D(A)$ pad a posterior by zeros outside $S_y$.
Given a continuation $E:S_y\to\D(O\times R_E)$ and arbitrary fixed
continuations $\mathbf B=(B_z)_{z\ne y}$ on the other branches, let
$J_{y,\mathbf B}(E)$ denote their compound with $P$.  Extend $E$ to the
other action rows arbitrarily: those rows have $P(y\mid a)=0$ and do not
affect the compound.  Tagged disjoint unions put all continuation records
in one finite alphabet when needed.

\begin{lemma}[Affine insertion]\label{mf:lem:insertion-affine}
There is a coefficient $c(q,P,y)>0$, independent of $\mathbf B$, such that
for every two continuations $E,G$ on $S_y$,
\begin{equation}
 W_q(J_{y,\mathbf B}(E))-W_q(J_{y,\mathbf B}(G))
 =c(q,P,y)[W_{r_y}(E)-W_{r_y}(G)].
 \label{mf:eq:unidentified-weight}
\end{equation}
This includes singleton posterior supports.
\end{lemma}

\begin{proof}
Support restriction and the one-branch form of Axiom~\textup{(A8)}
(Lemma~\ref{pt:lem:finite-branch-extension}\textup{(b)}) make insertion an
order embedding.  Its marked law is
\begin{equation}
 \eta_{J_{y,\mathbf B}(E)}
 =\beta_{\mathbf B}+m_y(\mathrm{id}_O,\iota_y)_\#\eta_E,
 \label{mf:eq:insertion-law}
\end{equation}
where $\beta_{\mathbf B}$, of mass $1-m_y$, is the contribution of the
other branches.  Insertion is therefore affine on marked laws.  The local
order is nonconstant because its two material anchors have values zero and
one.  Positive-affine uniqueness gives \eqref{mf:eq:unidentified-weight}
with an initially background-dependent positive coefficient.

Subtracting \eqref{mf:eq:insertion-law} for $E$ and $G$ removes
$\beta_{\mathbf B}$.  The resulting marked-law difference is the same for
every background.  Lemma~\ref{mf:lem:scaled-differences}, with $s=t=1$,
makes the numerical difference independent of the background as well.
Taking $E$ and $G$ to be the uninformative high and low payoffs, whose local
value difference is one, identifies the coefficients across backgrounds.
No information value has been used, and the argument applies when $S_y$ is
a singleton.
\end{proof}

For a payoff profile $g:A\to O$, let $\operatorname{Rev}(g)$ reveal the
realised action $a$ and pay $g(a)$.  Its marked law at $q$ is
\begin{equation}
 \eta_{\operatorname{Rev}(g)}
 =\sum_a q(a)\delta_{(g(a),\delta_a)}.
 \label{mf:eq:revelation-law}
\end{equation}
Write $E^-:=\operatorname{Rev}(o^-)$ and let $E^a$ pay $o^+$ at action
$a$ and $o^-$ at every other action, retaining full revelation.  Define
\begin{equation}
 \Delta_q(a):=W_q(E^a)-W_q(E^-),\qquad
 D_a:=\delta_{(o^+,\delta_a)}-\delta_{(o^-,\delta_a)}.
 \label{mf:eq:revealed-upgrade}
\end{equation}
Thus $\eta_{E^a}-\eta_{E^-}=q(a)D_a$.  Moreover, $\Delta_q(a)>0$:
view full revelation as the first stage and change only the payoff on the
reached branch $a$; the local high payoff is strictly preferred to the low
one, so the strict part of Axiom~\textup{(A8)} applies.

\begin{lemma}[A material unit at each prior]\label{mf:lem:material-units}
For every full-support prior $q$ there is a number $\gamma_q>0$ such that
\begin{equation}
 \Delta_q(a)=q(a)\gamma_q\quad\text{for every }a.
 \label{mf:eq:upgrade-unit}
\end{equation}
The unit is invariant under relabelling and equals one on singleton action
sets.  For every reached branch as above,
\begin{equation}
 c(q,P,y)\gamma_{r_y}=m_y\gamma_q.
 \label{mf:eq:unit-ratio}
\end{equation}
\end{lemma}

\begin{proof}
\emph{Equal increments per unit of probability.}
Let $v$ be the uniform prior on $A$.  Exact numerical invariance under
simultaneous action and record permutations makes $\Delta_v(a)$ the same
for every $a$.  Choose $0<\varepsilon<1$ so small that
$\varepsilon v(a)\le q(a)$ for every $a$.  The binary channel
\[
 P(1\mid a)=\frac{\varepsilon v(a)}{q(a)},\qquad
 P(0\mid a)=1-P(1\mid a)
\]
reaches posterior $v$ with probability $\varepsilon$ on record $1$.
Fix a low-payoff continuation after record $0$, and write $J$ for insertion
after record $1$, with coefficient $c>0$ from
Lemma~\ref{mf:lem:insertion-affine}.  At the outer prior $q$, the direct
upgrade has marked-law difference $q(a)D_a$, whereas the inserted upgrade
has difference $\varepsilon v(a)D_a$.  Consequently
Lemma~\ref{mf:lem:scaled-differences} gives
\[
 \varepsilon v(a)\Delta_q(a)
 =q(a)[W_q(J(E^a))-W_q(J(E^-))]
 =q(a)c\Delta_v(a).
\]
Both $v(a)$ and $\Delta_v(a)$ are independent of $a$.  Dividing by the
positive $\varepsilon v(a)q(a)$ shows that $\Delta_q(a)/q(a)$ is constant
in $a$.  Call it $\gamma_q$.  It is positive and, since $\sum_a q(a)=1$,
also satisfies $\gamma_q=\sum_a\Delta_q(a)$.  This expression and exact
relabel invariance of $W$ give relabel invariance of $\gamma$.  On a
singleton set the sole upgrade is the uninformative low-to-high change,
so $\gamma_q=1$.

\emph{The branch ratio.}
Choose $a\in S_y$ and compare the same high-payoff upgrade at the pure
posterior $\delta_a$.  Direct full revelation at $q$ gives the marked-law
difference $q(a)D_a$.  Full revelation on the branch support, inserted
after $y$, gives $m_y r_y(a)D_a$, with the posterior embedded by
$\iota_y$.  The scaled-difference lemma and affine insertion imply
\[
 m_y r_y(a)\Delta_q(a)
 =q(a)c(q,P,y)\Delta_{r_y}(a).
\]
Substitute \eqref{mf:eq:upgrade-unit} on both sides and cancel
$q(a)r_y(a)>0$.  This gives \eqref{mf:eq:unit-ratio}.
\end{proof}

The remaining calibration uses a partition with just two cells.  One cell
carries an arbitrary action distribution; the other contains one action.
Changing the material payoff on either cell can be evaluated before or
after the actions in that cell are merged.

\begin{lemma}[Exact reached-branch weight]\label{lf:lem:branch-weight}
Every material unit equals one.  Hence, for every full-support prior $q$,
every first-stage channel $P$, every reached $y$, all continuations $E,G$
on its posterior support, and every fixed background $\mathbf B$,
\begin{equation}
 W_q(J_{y,\mathbf B}(E))-W_q(J_{y,\mathbf B}(G))
 =m_y[W_{r_y}(E)-W_{r_y}(G)].
 \label{lf:eq:branch-weight}
\end{equation}
The formula also holds for arbitrary outer priors after restriction to
their positive supports.
\end{lemma}

\begin{proof}
Fix any full-support $r\in\D(A)$.  On $A\sqcup\{\ast\}$ put
\[
 Q(a)=\tfrac12r(a)\ (a\in A),\qquad Q(\ast)=\tfrac12,
\]
and let $P$ report the cell $A$ or $\{\ast\}$.  The cell lottery is
$p=(\tfrac12,\tfrac12)$; its two conditional action distributions are
$r$ and the singleton prior, up to the harmless support identification.
Write $c_1,c_2$ for the two insertion coefficients.

For any pair of cell payoffs $(g_1,g_2)$, consider the experiment that
reports the cell and pays its assigned payoff.  The action processor that
replaces each action by its cell maps this experiment to full revelation
of $p$ with the same two payoffs.  Conversely, from cell $1$ draw an action
according to $r$, and from cell $2$ choose $\ast$.  These processors are
Bayesian pushforwards in both directions: all rows within a cell have the
same visible output.  Axiom~\textup{(A7)} and the common numerical scale
therefore equate the values of the two experiments for every $(g_1,g_2)$.

Start with both cell payoffs equal to $o^-$ and upgrade only cell $i$.
Before merging, its value increment is $c_i$, since the local material
increment is one.  After merging, its value increment is
$\Delta_p(i)=\tfrac12\gamma_p$.  Thus
\begin{equation}
 c_1=c_2=\tfrac12\gamma_p.
 \label{mf:eq:two-cell-upgrades}
\end{equation}
The branch ratio \eqref{mf:eq:unit-ratio} on the singleton cell gives
$c_2=\tfrac12\gamma_Q$, hence $\gamma_Q=\gamma_p$.
On the other cell it gives $c_1\gamma_r=\tfrac12\gamma_Q$.
Substituting \eqref{mf:eq:two-cell-upgrades} and using $\gamma_p>0$
yields $\gamma_r=1$.  Since $r$ was arbitrary, every unit is one.
Equation~\eqref{mf:eq:unit-ratio} now gives $c(q,P,y)=m_y$; substitute
this in \eqref{mf:eq:unidentified-weight}.

For a boundary prior, delete its null action rows first.  Every reached
posterior lies in the remaining support, and the branch masses, conditional
values, and compound comparisons are unchanged by
Lemma~\ref{lf:lem:bookkeeping}.  This proves the final assertion.
\end{proof}

\subsection{Separating deterministic payoffs}

For a first-stage record channel $P:A\to\D(Y)$ and a payoff profile
$g:Y\to O$, let $C(P,g)$ be the compound that runs $P$, retains its record
$y$, and then pays $g(y)$ with no further record.  Write $C(P,o)$ when
$g\equiv o$, and define the value of the record at the low material anchor by
\begin{equation}
T_q(P):=W_q\bigl(C(P,o^-)\bigr).
\label{mf:eq:trace-value}
\end{equation}
As elsewhere, a boundary prior is evaluated after restriction to its positive
support.  Thus this definition applies to every prior, including a point mass.
No information formula for $T_q$ is assumed here.

\begin{lemma}[Deterministic payoff separation]
\label{lf:lem:deterministic-assembly}\label{lf:lem:branch-increment}
For every nonempty finite $A,Y$, every $q\in\D(A)$, every
$P:A\to\D(Y)$, and every $g:Y\to O$, put
$m_y:=\sum_aq(a)P(y\mid a)$.  Then
\begin{equation}
W_q(C(P,g))=T_q(P)+\sum_{y\in Y}m_yu(g(y)).
\label{lf:eq:deterministic-assembly}
\end{equation}
In particular, for every $o\in O$,
\begin{equation}
W_q(C(P,o))=T_q(P)+u(o).
\label{mf:eq:constant-payoff-shift}
\end{equation}
\end{lemma}

\begin{proof}
First let $q$ have full support.  For a profile $g$ and an outcome $o$, write
$g[y\leftarrow o]$ for the profile obtained by changing only the payoff at
$y$ to $o$.  If $m_y=0$, this change leaves the marked terminal law unchanged.
If $m_y>0$, the exact branch-weight identity in
Lemma~\ref{lf:lem:branch-weight}, applied to the two deterministic local
payoffs and with all other continuations fixed, gives
\begin{equation}
W_q\bigl(C(P,g[y\leftarrow o])\bigr)-W_q(C(P,g))
=m_y\bigl[u(o)-u(g(y))\bigr].
\label{lf:eq:branch-increment}
\end{equation}
Here the local deterministic-payoff values are $u(o)$ and $u(g(y))$ by
\eqref{lf:eq:payoff-normalisation}.  This remains true when the reached
posterior has singleton support.  The same identity holds when $m_y=0$,
because both sides are then zero.

Start from the constant profile $o^-$ and change the payoffs to $g(y)$ one
record at a time.  The arbitrary-background clause of
Lemma~\ref{lf:lem:branch-weight} permits each change after the preceding
ones have been made.  Summing \eqref{lf:eq:branch-increment} gives
\begin{equation}
W_q(C(P,g))-W_q(C(P,o^-))
=\sum_{y\in Y}m_y\bigl[u(g(y))-u(o^-)\bigr].
\label{lf:eq:telescope}
\end{equation}
Since $u(o^-)=0$, this is \eqref{lf:eq:deterministic-assembly}.
The proof includes a singleton action set: the only local values needed
are values of payoff lotteries.

For a boundary prior, restrict $q$ and the rows of $P$ to $\supp q$.
The branch probabilities $m_y$ do not change, and the support convention
preserves both compound values.  The full-support formula therefore proves
the stated identity for every $q$.  Finally, substituting $g\equiv o$ and
using $\sum_y m_y=1$ gives \eqref{mf:eq:constant-payoff-shift}.
\end{proof}

\subsection{Entropy reduction and identification}
\label{mf:sec:entropy}

For a prior $q\in\D(A)$, define the value of full revelation at the low
material payoff by
\[
 \mathcal H(q):=T_q(\Id_A).
\]
All values at boundary priors are read on their positive supports.  Relabelling
actions and records leaves $\mathcal H$ unchanged, as does adding or deleting
zero-probability actions.  Moreover, $\mathcal H(\delta_a)=0$: on a singleton
support, the action record is constant and can be erased or restored by
record processing.  Erasing the record of full revelation gives the constant
payoff $o^-$, whose value is zero.  Axiom~\textup{(A6)} therefore gives
$\mathcal H(q)\ge0$ for every $q$.

\begin{lemma}[Entropy reduction and grouping]
\label{mf:lem:entropy-reduction}
For every finite prior $q$ and record channel $P:A\to\D(Y)$, let $m_y$ be
the probability of record $y$ and $q_y$ its posterior on actions when $m_y>0$.
Then
\begin{equation}
 T_q(P)=\mathcal H(q)-\sum_{y:m_y>0}m_y\mathcal H(q_y).
 \label{mf:eq:entropy-reduction}
\end{equation}
For any finite probability vector $p=(p_i)_i$ and conditional probability
vectors $q^i$, on disjoint labelled sets,
\begin{equation}
 \mathcal H\bigl((p_iq^i)_i\bigr)
 =\mathcal H(p)+\sum_{i:p_i>0}p_i\mathcal H(q^i).
 \label{mf:eq:grouping}
\end{equation}
\end{lemma}

\begin{proof}
First suppose $q$ has full support.  Run $P$ and, after each reached record,
either reveal the action at payoff $o^-$ or deliver $o^-$ without further
information.  If every continuation reveals, the final marked law is
$\sum_aq(a)\delta_{(o^-,\delta_a)}$, so its value is $\mathcal H(q)$.
If every continuation is silent, the value is $T_q(P)$.  Change the reached
continuations one at a time.  Lemma~\ref{lf:lem:branch-weight} gives the
increment $m_y\mathcal H(q_y)$, because the silent local continuation has
value zero; singleton posterior supports are included.  Summing these
increments proves \eqref{mf:eq:entropy-reduction}.  A null record contributes
nothing, and restriction to the positive support proves the formula for
arbitrary $q$.

For grouping, first take all $p_i>0$ and all conditional vectors on their
positive supports.  Write $q(i,a)=p_iq^i(a)$ and let $C$ report only the cell
$i$.  Collapsing the action $(i,a)$ to $i$ transforms this experiment into
full revelation at prior $p$.  Conversely, after action $i$, draw $a$ from
$q^i$ and output $(i,a)$; the record remains $i$.  These two action processors
satisfy the exact joint-law equation of Axiom~\textup{(A7)} in their respective
directions.  Hence $T_q(C)=\mathcal H(p)$ on the common material scale.
The posterior after $i$ is $q^i$ on that cell.  Substitution in
\eqref{mf:eq:entropy-reduction} gives \eqref{mf:eq:grouping}.  Deleting and
reinserting zero-probability actions or cells gives its stated general form.
\end{proof}

\begin{lemma}[Shannon identification]
\label{mf:lem:entropy-identification}
There is a single $\lambda>0$, independent of all finite alphabets and
priors, such that
\begin{equation}
 \mathcal H(q)=\lambda\Sh(q),
 \qquad
 T_q(P)=\lambda I_{qP}(A;Y)
 \label{mf:eq:trace-MI}
\end{equation}
for every finite prior $q$ and record channel $P:A\to\D(Y)$.
\end{lemma}

\begin{proof}
\emph{Binary continuity from material comparisons.}
Put $q_t=(t,1-t)$ and $h(t):=\mathcal H(q_t)$.  We first prove continuity
on $(0,1)$.  Let $E$ be binary full revelation at payoff $o^-$, and let $M$
be the action-independent material lottery assigning probability $1/2$ to
each of $o^+$ and $o^-$.  Thus $W_{q_t}(E)=h(t)$ and $W_{q_t}(M)=1/2$.
Fix any real threshold $c$.  Choose $\varepsilon\in(0,1)$ small enough that
\[
 b:=\varepsilon c+(1-\varepsilon)/2\in(0,1),
\]
and let $B_b$ be the action-independent material lottery paying $o^+$ with
probability $b$ and $o^-$ otherwise.  The public mixture
$E_c:=\varepsilon E\oplus(1-\varepsilon)M$ and the benchmark $B_b$ are
fixed experiments as $t$ varies.  Affinity and material normalisation give
\[
 W_{q_t}(E_c)-W_{q_t}(B_b)
 =\varepsilon\bigl(h(t)-c\bigr).
\]
Consequently the sets $\{t:h(t)\ge c\}$ and $\{t:h(t)\le c\}$ are
comparison sets in the fixed two-block environment containing $E_c$ and
$B_b$.  The block-supported priors vary continuously with $t$, so
Axiom~\textup{(A2)} makes both sets closed relative to $(0,1)$.
This holds for every $c\in\R$, proving interior continuity.

For completeness, grouping also supplies continuity at the endpoints.
For $0<t<1/4$, evaluate the vector $(t,1/2-t,1/2)$ by first merging its
first two coordinates and then by isolating its first coordinate.  Equation
\eqref{mf:eq:grouping} gives
\[
 h(t)=g(t)+\tfrac12h(2t),
 \qquad
 g(t):=h(\tfrac12)-(1-t)
 h\!\left(\frac{1/2-t}{1-t}\right).
\]
Interior continuity makes $g$ continuous on $[0,1/4]$ after setting $g(0)=0$.
Choose $n$ so that $2^nt\in[1/4,1/2)$ and iterate:
\[
 h(t)=\sum_{j=0}^{n-1}2^{-j}g(2^jt)+2^{-n}h(2^nt).
\]
Both $g$ on $[0,1/4]$ and $h$ on $[1/4,1/2]$ are bounded.
As $t\downarrow0$, the terminal term tends to zero.  For the sum, any fixed
finite initial segment tends to zero and the remaining tail is uniformly
bounded by a geometric tail.  Thus $h(t)\to0=h(0)$.  Symmetry gives
continuity at one.

\emph{Identification and strict positivity.}
We use the finite Faddeev characterisation: a nonnegative family on finite
probability vectors, symmetric and unchanged by adjoining zeros, vanishing
at point masses, satisfying the grouping identity, and having a continuous
binary restriction is a common nonnegative multiple of Shannon entropy
\cite{Faddeev1956}.  This is the binary-continuity form of the theorem;
continuity on every closed finite simplex need not be assumed separately.
The preceding paragraphs establish binary continuity, and
Lemma~\ref{mf:lem:entropy-reduction} supplies grouping.  The remaining
properties were established directly after the definition of $\mathcal H$.
Hence $\mathcal H(q)=\lambda\Sh(q)$ for some common $\lambda\ge0$.

To exclude $\lambda=0$, let $o_\ast$ witness Axiom~\textup{(A4)}.
Lemma~\ref{lem:relevance-bridge} identifies its strict comparison with fair
binary full revelation versus silence at payoff $o_\ast$.  By deterministic
payoff separation (Lemma~\ref{lf:lem:deterministic-assembly}), these two
values are respectively
\[
 u(o_\ast)+\mathcal H(\tfrac12,\tfrac12)
 \quad\text{and}\quad u(o_\ast).
\]
Therefore $\mathcal H(1/2,1/2)>0$ and $\lambda>0$.
Finally, substitution in \eqref{mf:eq:entropy-reduction} and the finite
posterior-entropy identity give
\[
 T_q(P)
 =\lambda\left[\Sh(q)-\sum_{y:m_y>0}m_y\Sh(q_y)\right]
 =\lambda I_{qP}(A;Y).
\]
Support restriction covers boundary priors, and the multiplier is the same
for every action and record alphabet.
\end{proof}

\subsection{Proof of the theorem}

\begin{corollary}[Value of a deterministic payoff continuation]
\label{mf:cor:deterministic-value}
Let $\lambda>0$ be the coefficient in
Lemma~\ref{mf:lem:entropy-identification}.  For every finite first-stage
record channel $P:A\to\D(Y)$, every $q\in\D(A)$, and every $g:Y\to O$,
\[
W_q(C(P,g))=\lambda I_{qP}(A;Y)+\sum_{y\in Y}m_yu(g(y)),
\qquad m_y:=\sum_aq(a)P(y\mid a).
\]
\end{corollary}

\begin{proof}
Substitute $T_q(P)=\lambda I_{qP}(A;Y)$ into
\eqref{lf:eq:deterministic-assembly}.
\end{proof}

\begin{proof}[Proof of Theorem~\ref{thm:main}]
Assume Axioms~\textup{(A1)--(A8)}.  The material normalisation supplies the
nonconstant $u$ with $u(o^-)=0$ and $u(o^+)=1$, and
Lemma~\ref{mf:lem:entropy-identification} supplies one $\lambda>0$ such that
$T_q(P)=\lambda I_{qP}(A;Y)$ for every finite record experiment.

Fix a full-support $q\in\D(A)$ and an arbitrary
$K:A\to\D(O\times R)$.  Factor $K$ through the first-stage record
$Y=O\times R$ by setting
\[
P_K(y\mid a):=K(o,r\mid a),\qquad y=(o,r),
\]
and then paying $g(o,r):=o$ with no further record.  Write
$K^{\mathrm{seq}}:=C(P_K,g)$ for this compound.  The final record identifies
$(o,r)$; its posterior over actions and its realised payoff are exactly those
under $K$.  Hence the two experiments have the same marked terminal law and
the same $W_q$-value, by Lemma~\ref{lf:lem:terminal-affine}.
Deterministic payoff separation gives
\[
W_q(K)=T_q(P_K)+\sum_{o,r}p_{qK}(o,r)u(o).
\]
Since $I_{qP_K}(A;Y)=I_{qK}(A;O,R)$, entropy identification now yields
\begin{equation}
W_q(K)=\E_{qK}[u(O)]+\lambda I_{qK}(A;O,R)=:V(q,K).
\label{lf:eq:full-support-value}
\end{equation}
This also covers a singleton action set.  In that case every posterior is
the same point mass, information is zero, and terminal-law sufficiency and
\eqref{lf:eq:payoff-normalisation} give the same formula directly.

For an arbitrary prior $q$, put $S:=\supp q$ and let $(q_S,K_S)$ denote
the restricted experiment.  Support transport and the support convention
for $W$ give
\begin{equation}
(q,K)\sim(q_S,K_S),\qquad
W_q(K)=W_{q_S}(K_S)=V(q_S,K_S)=V(q,K).
\label{lf:eq:support-value}
\end{equation}
The last equality follows directly from the finite sums for expected utility
and mutual information: omitted actions have zero probability.  Thus
\eqref{lf:eq:full-support-value} extends to every prior.

For any two finite experiments $(q,K)$ and $(p,L)$, restrict both priors to
their supports.  The common material scale
\eqref{lf:eq:common-scale} represents their comparison there, and
\eqref{lf:eq:support-value} transports both values and the comparison back.
Consequently,
\begin{equation}
(q,K)\succeq(p,L)
\quad\Longleftrightarrow\quad
V(q,K)\ge V(p,L),
\label{lf:eq:pair-value}
\end{equation}
with the same $u$ and $\lambda$ for every finite alphabet.  The duplication
clause of Axiom~\textup{(A5)} gives the required within-channel
representation.  Its irrelevant-block clause applies
\eqref{lf:eq:pair-value} to any two block-supported priors in any finite
block environment, proving the moreover clause with these same witnesses.

Conversely, suppose a nonconstant $u:O\to\R$ and $\lambda>0$ represent
the orders in every finite environment by $V$.
Completeness and transitivity of the numerical order give
Axiom~\textup{(A1)}.  On fixed finite alphabets, expected utility and mutual
information are jointly continuous in $(q,K)$; for the latter, use the finite
entropy formula and $0\log0=0$.  The difference $V(q,K)-V(p,K)$ is therefore continuous in $(K,q,p)$,
so its nonnegative set is closed.  This proves Axiom~\textup{(A2)}.

Under a prior supported on block $i$ of a finite block channel, the block tag
is constant and the payoff law and information agree with those of the
untagged experiment.  Thus
\[
V\left(q_i^i,\bigsqcup_{k\in I}K_k\right)=V(q_i,K_i).
\]
For duplicated channels this proves both directions of the duplication
clause of Axiom~\textup{(A5)}.  For arbitrary distinct blocks it proves
both directions of the irrelevant-block clause and the theorem's moreover
clause, using the same $u$ and $\lambda$.

A payoff-preserving record processor leaves the payoff distribution
unchanged and, by the Markov chain $A-(O,R)-(O,R')$, weakly lowers mutual
information.  This proves Axiom~\textup{(A6)}.
For action processing, the joint law
\[
q(a)K(o,r\mid a)S(b\mid a)
\]
has Markov chain $(O,R)-A-B$.  Its $(B,O,R)$ marginal is precisely the
Bayesian pushforward experiment in \eqref{eq:action-processing}, and its
$(O,R)$ marginal is unchanged.  Hence
$I(B;O,R)\le I(A;O,R)$ proves Axiom~\textup{(A7)}.
Every row completed at a $b$ with $(qS)(b)=0$ has zero joint mass, so the
argument applies to every completion allowed by the axiom.

For Axiom~\textup{(A8)}, take a binary first-stage channel $P$, a prior $q$,
and the continuation triple $(K,L,M)$ covered by the axiom.  Write
$\alpha:=\sum_aq(a)P(1\mid a)>0$ and let $q_1$ be the posterior after the
first record.  The expected-utility and mutual-information chain rules give
\begin{align*}
&V\bigl(q,P\triangleright(K,M)\bigr)
 -V\bigl(q,P\triangleright(L,M)\bigr)\\
&\qquad=\alpha\bigl[V(q_1,K)-V(q_1,L)\bigr].
\end{align*}
The first-stage information term and the common continuation after record
$2$ cancel.  Since $\alpha>0$, the aggregate comparison is weakly positive
exactly when the conditional comparison is weakly positive, and the same
equivalence holds for strict positivity.  This proves Axiom~\textup{(A8)}.

Finally, nonconstancy of $u$ supplies $o^+,o^-$ with $u(o^+)>u(o^-)$.
In the two-action channel of Axiom~\textup{(A3)}, the two pure lotteries have
these utilities and zero mutual information, establishing material
relevance.  In the four-action channel of Axiom~\textup{(A4)}, both tested
lotteries pay the same outcome surely, while their information values are
$\log2$ and zero.  The required strict comparison follows from $\lambda>0$.
All eight axioms, and the same-witness block clause, have therefore been
verified.
\end{proof}
\clearpage
\section{Further results and choice implications}\label{app:auxiliary-results}

This appendix proves the remaining results, in the order in which they
appear in the text: uniqueness and signs, independence of the axioms,
choice and measurement, the moral-hazard application, and interpretation.
None is used in Appendix~\ref{app:main-proof}.
Unless a result specifies weaker axioms or a different sign, we use
$V(q,K)=E_{qK}+\lambda\I_{qK}$ with $\lambda>0$, where
$E_{qK}:=\E_{qK}[u(O)]$ and $\I_{qK}:=\I_{qK}(A;O,R)$.
The posterior on actions after $(o,r)$ and its distribution are
\[
 \pi^{qK}_{or}(a):=\frac{q(a)K(o,r\mid a)}{p_{qK}(o,r)},
 \qquad
 \mu_{qK}:=\sum_{o,r}p_{qK}(o,r)\,\delta_{\pi^{qK}_{or}},
\]
with zero-probability outputs omitted.

\subsection{Consequences of the representation}
\label{app:representation-consequences}

\begin{corollary}[No payoff--trace complementarity]\label{cor:matching}
Suppose Axioms \textup{(A1)--(A8)} hold.  Fix $A$ and $q\in\D(A)$.  If channels $K$ and $L$ satisfy
$E_{qK}=E_{qL}$ and $\mu_{qK}=\mu_{qL}$, then $(q,K)\sim(q,L)$.
\end{corollary}

\begin{proof}
The posterior-entropy identity
$\I_{qK}(A;O,R)=\Sh(q)-\int\Sh(\pi)\,d\mu_{qK}(\pi)$
gives equal information values.  The common-scale representation in
Theorem~\ref{thm:main} then gives indifference.
\end{proof}

\begin{proof}[Proof of Proposition~\ref{prop:uniqueness}]
\emph{Range and ordinal uniqueness.}
The range of $V$ is $[\underline u,\infty)$.
The lower bound follows from nonnegative information.  To attain every value
above it, choose $o_*$ with $u(o_*)=\underline u$, use the uniform prior on
$n$ actions, and pay $o_*$ surely.  Reveal the action with probability $t$
and otherwise report a distinct erasure symbol.  The resulting value is
$\underline u+\lambda t\log n$.  Varying $t\in[0,1]$ and $n\ge2$ fills
the claimed range.

If $W$ is common-scale, define $\varphi(v):=W(q,K)$ using any experiment
with $V(q,K)=v$.  Two experiments with the same $V$-value are indifferent
in their two-block environment, so their $W$-values agree; strict comparisons
likewise make $\varphi$ strictly increasing.  Thus $W=\varphi\circ V$.
Conversely, a single strictly increasing transformation preserves every
comparison represented by $V$.  This proves part~\textup{(i)}.

\emph{Affine uniqueness.}
The expected-utility and information chain rules give
\eqref{eq:branch-additive-index} for $V$: the common first-stage information
term cancels, leaving the probability-weighted sum of continuation-value
differences.  Every positive affine transformation of $V$
is therefore common-scale and branch-additive.

For the converse, write $W=\varphi\circ V$ and set
\[
 \psi(x):=\varphi(\underline u+x)-\varphi(\underline u),\qquad x\ge0.
\]
At the uniform prior on $n$ actions, let $E_x$ be the preceding erasure
experiment with value $\underline u+x$, where $0\le x\le\lambda\log n$,
and let $E_0$ be silence at the same payoff.  After an action-independent
coin with probabilities $\theta$ and $1-\theta$, compare continuations
$(E_x,E_0)$ with $(E_0,E_0)$.  Their $V$-values are
$\underline u+\theta x$ and $\underline u$, respectively.  Branch
additivity of $W$ therefore gives
\[
 \psi(\theta x)=\theta\psi(x),\qquad 0\le\theta\le1.
\]
The continuation records can use the common alphabet consisting of the
action labels and the erasure symbol; the formula also covers zero coin
weights by omitting unreached branches.  Taking $x=\lambda\log n$ makes
$\psi$ linear on $[0,\lambda\log n]$.  These nested, unbounded intervals
have a common positive-length overlap, so their slopes agree.  Strict
increase makes the common slope $\alpha$ positive.  Hence
$W=\alpha V+\beta$, proving part~\textup{(ii)}.

For the parameter claim, a second trace-tempered index $\widetilde V$ is
common-scale and branch-additive by the same chain rules.  The affine
identity, evaluated first at deterministic payoffs and then at a
constant-payoff informative experiment, gives
$\widetilde u=\alpha u+\beta$ and $\widetilde\lambda=\alpha\lambda$.
\end{proof}

\begin{proof}[Proof of Corollary~\ref{cor:signduality}]
Part~\textup{(i)} is Theorem~\ref{thm:main}.
For part~\textup{(ii)}, applying each reversed processing axiom to an
invertible relabelling and its inverse gives comparisons in both directions.
With A5, the block identifications of Lemma~\ref{lf:lem:bookkeeping}
therefore still apply.  In particular, swapping the two block tags identifies
a reversed canonical comparison with the reverse comparison inside the
original block.  Now reverse the primitive order in every environment.
The reversed family
satisfies \textup{(A1)}, \textup{(A2)}, \textup{(A5)}, and \textup{(A8)},
as well as the forward versions of \textup{(A4)}, \textup{(A6)}, and
\textup{(A7)}.  It satisfies \textup{(A3)} with the two witnessing outcomes
and their action labels interchanged.  Theorem~\ref{thm:main} represents it by
$\E[\widehat u(O)]+\widehat\lambda\I(A;O,R)$ with
$\widehat\lambda>0$.  Negating this index represents the original family,
with $u=-\widehat u$ and $\lambda=-\widehat\lambda<0$.

For part~\textup{(iii)}, the indifference versions of processing imply
their weak forward versions.  The payoff-lottery part of
Lemma~\ref{lf:lem:material-scale}, which does not require trace relevance,
therefore supplies a nonconstant utility $u$.  By $(\mathrm{A}7^0)$, collapsing
the action alphabet makes every experiment indifferent to a singleton-action
experiment with the same output law.  By $(\mathrm{A}6^0)$, deleting its
record leaves precisely its payoff lottery.  Its value is consequently
$\E[u(O)]$ by \eqref{lf:eq:payoff-normalisation}.  A5 gives the same
representation for within-channel and finite block comparisons.
Conversely, expected utility satisfies all three indifference clauses.
The positive and negative formulae satisfy their respective processing and
relevance clauses by data processing, and their sure-thing clauses by the
two chain rules.
\end{proof}

\subsection{Independence of the axioms}\label{app:axiom-independence}

\begin{proof}[Proof of Proposition~\ref{prop:independence}]
Fix outcomes $o^-,o^+$ and a utility $u$ with $u(o^-)=0$, $u(o^+)=1$.
Write $Z=(O,R)$, $U(q,K)=\E_{qK}[u(O)]$, and
$I(q,K)=\I_{qK}(A;Z)$.  Each construction below defines the primitive
order \emph{within every channel}; cross-channel statements are its canonical
two-block comparisons.

We use three common facts.  The coordinates $U$ and $I$ are continuous,
unchanged by block embedding, and weakly decrease under either processing
operation ($U$ is unchanged).  Under a replacement at a reached branch of
mass $m$, their aggregate differences equal $m$ times their conditional
differences.  These identities hold on boundary priors and when the other
branch is unreached.  Finally, the material test has coordinate difference
$(1,0)$ and the trace test has difference $(0,\log2)$.

\emph{Without \textup{(A1)}: incompleteness.}
Use the coordinatewise order on $(U,I)$.  It is transitive and has a closed
graph.  Block invariance gives A5, processing gives A6--A7, and multiplication
of both coordinate differences by the same $m>0$ gives the weak biconditional
in A8, including its strict part.  Both relevance tests are strict.
Completeness fails within the fully revealing two-action channel paying
$o^+$ at the first action and $o^-$ at the second: the pure high-payoff
lottery has coordinates $(1,0)$ and the fair lottery has
$(1/2,\log2)$, so neither dominates the other.

\emph{Without \textup{(A2)}: lexicographic priorities.}
Order $(U,I)$ lexicographically, with payoff first.  This is a weak order;
the preceding identities verify A3--A8, since multiplication by $m>0$
preserves lexicographic comparisons.  For a failure of closedness in one
fixed channel, reveal three actions fully and pay $o^+$ only at action $2$.
Put
\[
 p=(1/4,1/2,1/4),\qquad q=(1/2,1/2,0),\qquad
 q_n=(1/2-\varepsilon_n,1/2+\varepsilon_n,0),
\]
where $0<\varepsilon_n<1/2$ and $\varepsilon_n\downarrow0$.
Every $q_n$ has higher payoff than $p$, whereas $p$ and $q$ tie in payoff
and have information $\tfrac32\log2$ and $\log2$.  Thus
$q_n\succ_Kp$ for every $n$ but $q\not\succeq_Kp$.

\emph{Without \textup{(A3)} or \textup{(A4)}.}
For A3 use the score $I$.  All pure lotteries have value zero, while the
trace test remains strict.  For A4 use the score $U$.  Material relevance
holds, but the trace-test lotteries are indifferent.  The common facts
verify the other six axioms in each case, including A8 whenever a branch
is reached.

\emph{Without \textup{(A5)}: a bonus for full support.}
For each action alphabet $A$, use
\[
 V(q,K)=U(q,K)+I(q,K)+\prod_{a\in A}q(a).
\]
The added term is continuous, so A1--A2 hold.  It vanishes in both relevance
tests, since each lottery leaves some action unused.  It also vanishes on
every block-supported lottery in a canonical two-block comparison: the
other block is nonempty and has probability zero.  Thus all comparisons in
A6--A8, including the conditional comparisons in A8, agree with those of $U+I$.
Now take a two-action channel paying the same $o$ surely and leaving a
silent record.  The fair lottery has value $u(o)+1/4$, whereas a pure lottery
has value $u(o)$.  In the duplicated channel, their block-supported copies
both have value $u(o)$ and are indifferent.  The duplication clause of A5 fails.

\emph{Without \textup{(A6)}: a cost of conditional record noise.}
For $\varepsilon>0$, use
\begin{equation}
 V(q,K)=U(q,K)+I(q,K)-\varepsilon\Sh_{qK}(Z\mid A).
 \label{eq:no-record-dp}
\end{equation}
Here $\Sh_{qK}(Z\mid A)=\sum_aq(a)\Sh(K(\cdot\mid a))$ is continuous
and block invariant.  In a compound it equals
$\Sh(Y\mid A)+\sum_y m_y\Sh(O,R\mid A,Y=y)$.
Thus its replacement differences, like those of $U$ and $I$, are exactly
$m$ times the conditional differences: A8 holds, including strict
comparisons and singleton posterior supports.  A1, A2, and A5 follow
immediately; A3 has zero conditional entropy, and the two A4 scores differ
by $(1+\varepsilon)\log2$.
For action processing $A\to B$, the common output marginal and the Markov
chain $Z-A-B$ give the original-minus-processed difference
\[
 [I(A;Z)-I(B;Z)]
 +\varepsilon[\Sh(Z\mid B)-\Sh(Z\mid A)]
 =(1+\varepsilon)I(A;Z\mid B)\ge0.
\]
So A7 holds for every Bayesian completion.  But at a constant payoff, a fair
record independent of the action has value $u(o)-\varepsilon\log2$;
erasing it gives $u(o)$, violating A6.

\emph{Without \textup{(A7)}: Gini posterior reduction.}
Let $G(q)=1-\sum_aq(a)^2$ and use
\[
 J_G(q,K)=G(q)-\int G(\pi)\,d\mu_{qK}(\pi),\qquad V=U+J_G.
\]
This is block invariant.  It is also continuous at null outputs: writing
$x_{az}=q(a)K(z\mid a)$ and $p_z=\sum_a x_{az}$, its formula is
\[
 J_G(q,K)=\sum_{z:p_z>0}\frac{\sum_a x_{az}^2}{p_z}-\sum_aq(a)^2,
 \qquad 0\le\frac{\sum_a x_{az}^2}{p_z}\le p_z.
\]
Concavity of $G$ gives A6, because a garbled posterior is a conditional
average of the original posteriors.  Adding and subtracting
$\sum_y m_yG(q_y)$ gives the exact chain identity
\[
 J_G(q,P\triangleright(K_y)_y)
 =J_G(q,P)+\sum_{y:m_y>0}m_yJ_G(q_y,K_y),
\]
where $J_G(q,P)$ uses posteriors after the first-stage record alone.
Hence A8 holds, including its strict and boundary cases.  A1, A2, and A5
follow, pure lotteries have $J_G=0$, and the A4 values of $J_G$ are
$1/2$ and zero, verifying A3--A4.
To violate A7, take $q=(1/4,1/4,1/2)$, keep the payoff constant, and let the
record distinguish $\{a_1,a_2\}$ from $\{a_3\}$.  Before merging the
identical first two action rows, $J_G=5/8-(1/2)(1/2)=3/8$.
After the merge the prior is $(1/2,1/2)$ and the record reveals it fully,
so $J_G=1/2$.  Action processing has increased value.

\emph{Without \textup{(A8)}: satiation.}
For any $c>0$, use $V=U+\min\{I,c\}$.  Continuity, block invariance,
processing monotonicity, and both relevance tests give A1--A7.
For a violation of A8, take the uniform prior $q$ on $N$ actions with
$\log N>2c$, keep the payoff constant, and let the first-stage channel $P$
report a fair coin independent of the action.  Let $K$ produce a constant record and $L$
reveal the action, using the common record alphabet of action labels and a
silence symbol.  Compare $P\triangleright(K,L)$ and $P\triangleright(L,L)$.
After record $1$, the posterior is still $q$, so the conditional capped
information values are $0$ and $c$.  The two compounds have information
$\tfrac12\log N$ and $\log N$, respectively, and hence both have capped
value $c$.  Thus the aggregate weak comparison holds by indifference while
its conditional comparison fails.  This violates the binary biconditional
in A8 and completes the eight constructions.
\end{proof}

\subsection{Choice implications}\label{app:choice-proofs}

We prove the trace test and then the optimisation, record-comparison, and
measurement results used in Section~\ref{sec:implications}.  Write
$Z:=O\times R$, $K_a:=K(\cdot\mid a)$, and
$\bar u(a):=\E_{K_a}[u(O)]$; the output law under $q$ is $p_{qK}$.
Except in the explicitly axiomatic lemma below, we use the representation
with $\lambda>0$.  Write $b_0>1$ for the paper's fixed logarithm base.
Divergences use this base and
equal $+\infty$ when their first argument is not supported by their second.

\subsubsection{The trace test}\label{app:trace-test-proof}

\begin{proof}[Proof of Proposition~\ref{prop:signature}]
Both lotteries induce the uniform distribution on
$\{(o_0,r_1),(o_0,r_2)\}$.  Under $q'$ the record identifies the action, so
information is $\log2$; under $q''$ the two active rows coincide, so
information is zero.  The common payoff is $o_0$, giving the value gap
$\lambda\log2$.  In $K^\circ$ all rows coincide, so every lottery has
information zero and value $u(o_0)$, proving part~\textup{(ii)}.
\end{proof}

A channel-blind index $F(p_{qK},q)$ assigns each of the two lotteries the
same value in both treatments: neither its lottery argument nor its output
distribution changes.  It therefore cannot reproduce strict preference
in the first treatment and indifference in the second.  If $F(p,\cdot)$
is also invariant under action permutations, the first treatment alone
rejects it, because $q'$ and $q''$ are permutations with the same output law.
The gap in the first treatment has the sign of $\lambda$; the second
treatment gives indifference for positive, negative, and zero $\lambda$.

\subsubsection{Deliberate randomisation}
\label{app:randomisation-proofs}

\begin{lemma}[Qualitative trace preferences]\label{lem:qualitative-trace}
Suppose Axioms~\textup{(A1)} and~\textup{(A5)--(A7)} hold.  Let
$K:A\to\D(O\times R)$ pay the same $o\in O$ surely under every action.
Then $q\succeq_K\delta_a$ for every $q\in\D(A)$ and $a\in A$.
If action and record are independent under $qK$, then
$q\sim_K\delta_a$ for every $a$.  If Axiom~\textup{(A4)} also holds,
there are such a $K$ and $q$ with $q\succ_K\delta_a$ for every $a$.
\end{lemma}

\begin{proof}
Let $E_o$ be the one-action experiment paying $o$ with a silent record.
If action and record are independent under $qK$, delete null actions by
Lemma~\ref{lf:lem:bookkeeping}\textup{(b)}.  The remaining rows coincide.
Collapsing these actions to one, or drawing them back according to $q$,
are inverse action processings of the joint law, so Axiom~\textup{(A7)}
gives indifference.  Erasing or independently regenerating the common
record gives indifference to $E_o$ by Axiom~\textup{(A6)}.  In particular,
every pure lottery in $K$ is indifferent to $E_o$.

For arbitrary $q$, erase its record.  The result is indifferent to $E_o$,
so $(q,K)\succeq(\delta_\ast,E_o)\sim(\delta_a,K)$.  Replacement and duplication
(Lemma~\ref{lf:lem:bookkeeping}\textup{(a)} and Axiom~\textup{(A5)})
give the within-channel comparison.  Finally, the noisy coin in
Axiom~\textup{(A4)} has independent action and record and is therefore
indifferent to every pure action.  The strictly preferred revealing coin
proves the last assertion.
\end{proof}

\begin{samepage}
\begin{proposition}[Optimal mixing]\label{prop:binary}
Let $A=\{a,b\}$, $K_a\ne K_b$, and
$\Delta:=\bar u(a)-\bar u(b)\ge0$.  Let $q_x(b)=x$ and
$p_x:=(1-x)K_a+xK_b$.  For $0<x<1$, put
\[
g(x):=\KL(K_b\|p_x)-\KL(K_a\|p_x).
\]
\begin{enumerate}[label=\textup{(\roman*)},leftmargin=*]
\item $g$ is continuous and strictly decreasing, with extended-real limits
$g(0^+)=\KL(K_b\|K_a)$ and $g(1^-)=-\KL(K_a\|K_b)$.
The pure lottery $\delta_a$ is optimal exactly when
$\Delta\ge\lambda\KL(K_b\|K_a)$.  Below this threshold the unique
optimum is interior, with mass $x^*$ on $b$ determined by
$\Delta=\lambda g(x^*)$.
\item The optimal mass $x^*(\Delta,\lambda)$ is nonincreasing in $\Delta$
and nondecreasing in $\lambda$.  In the interior it decreases strictly
with $\Delta$ and, when $\Delta>0$, increases strictly with $\lambda$.
At $\Delta=0$ it is the capacity-achieving input of the two-row channel,
independently of $\lambda$.
\end{enumerate}
\end{proposition}
\end{samepage}

\begin{proof}
Write
\[
f(x):=V(q_x,K)=\bar u(a)-x\Delta+\lambda I(x),\qquad
I(x):=\Sh(p_x)-(1-x)\Sh(K_a)-x\Sh(K_b).
\]
On $(0,1)$, every consequence reached by either row has positive
$p_x$-mass.  Differentiating on this union of supports gives
\[
I'(x)=g(x),\qquad
g'(x)=-\frac1{\ln b_0}\sum_{z: p_x(z)>0}
\frac{(K_b(z)-K_a(z))^2}{p_x(z)}<0.
\]
At zero, supported terms converge
to the corresponding divergence terms; any $K_b(z)>0=K_a(z)$ contributes
$+\infty$.  This gives $g(0^+)$; exchanging the rows gives $g(1^-)$.
Thus $f$ is strictly concave.  Its right derivative at zero is nonpositive
exactly at the stated threshold.  Otherwise that derivative is positive,
whereas $f'(1^-)=-\Delta-\lambda\KL(K_a\|K_b)<0$, giving one interior
root.  There $x^*=g^{-1}(\Delta/\lambda)$, which proves all the comparative
statics; at the boundary $x^*=0$.  When $\Delta=0$, maximising $f$ is
exactly maximising channel information.
\end{proof}

\paragraph{General optimality.}
The finite entropy formula gives
\begin{equation}\label{eq:choice-value}
\begin{aligned}
V(q,K)
&=\lambda\Sh(p_{qK})+
  \sum_aq(a)\bigl[\bar u(a)-\lambda\Sh(K_a)\bigr]\\
&=\sum_{a\in\supp q}q(a)
  \bigl[\bar u(a)+\lambda\KL(K_a\|p_{qK})\bigr].
\end{aligned}
\end{equation}
The first line proves continuity and concavity in $q$, since entropy is
continuous and concave and the other term is linear.  Compactness of
$\D(A)$ therefore gives an optimum.  For $p\in\D(Z)$, define
$\Phi_a(p):=\bar u(a)+\lambda\KL(K_a\|p)$.
Along $q_t=(1-t)q+t\delta_a$,
\begin{equation}\label{eq:choice-direction}
\left.\frac{d}{dt}V(q_t,K)\right|_{t=0^+}
=\Phi_a(p_{qK})-V(q,K).
\end{equation}
When $\supp K_a\subseteq\supp p_{qK}$ this follows by differentiating
the first line of \eqref{eq:choice-value}; the constant derivative terms
cancel because both output distributions sum to one.  Otherwise, if $d>0$
is $K_a$'s mass outside that support, those new coordinates contribute
$\lambda d\,t\log(1/t)+O(t)$, so the derivative is $+\infty$, as is the
right-hand side.

We will also use the following exact identity.  If $p$ is positive on the
support of every row, expansion of the logarithms gives, for every $r$,
\begin{equation}\label{eq:choice-reference}
V(r,K)=\sum_a r(a)\Phi_a(p)-\lambda\KL(p_{rK}\|p).
\end{equation}

\begin{proof}[Proof of Proposition~\ref{prop:optimal-choice}]
At an optimum $q^*$, every derivative in \eqref{eq:choice-direction} is
nonpositive.  Hence $\Phi_a(p_{q^*K})\le V(q^*,K)$ for every $a$.
Their $q^*$-weighted average equals $V(q^*,K)$ by
\eqref{eq:choice-value}, so equality holds on $\supp q^*$.
Conversely, the stated conditions imply $c=V(q^*,K)$ and make all row
divergences finite.  Apply \eqref{eq:choice-reference} with $p=p_{q^*K}$:
for every $r$,
$V(r,K)\le\sum_a r(a)\Phi_a(p)\le c=V(q^*,K)$.
Thus the conditions are sufficient as well as necessary.  At constant
material utility they reduce to the channel-capacity condition.
\end{proof}

\begin{samepage}
\begin{proposition}[Support and pure choice]\label{cor:support-pure}
Let $Z_K:=\{z:K_a(z)>0\text{ for some }a\}$.
\begin{enumerate}[label=\textup{(\roman*)},leftmargin=*]
\item Every optimal $q^*$ has $p_{q^*K}(z)>0$ for all $z\in Z_K$.
Consequently, if each action $a$ has a private consequence $z_a$ with
$K_a(z_a)>0$ and $K_b(z_a)=0$ for all $b\ne a$, every optimum has full
support.
\item The pure lottery $\delta_a$ is optimal if and only if
\[
\bar u(b)+\lambda\KL(K_b\|K_a)\le\bar u(a)
\qquad\text{for every }b,
\]
and uniquely optimal if the inequalities are strict for $b\ne a$.
If a rival reaches any consequence outside $\supp K_a$, then $\delta_a$
is not optimal, however large its payoff advantage.
\end{enumerate}
\end{proposition}
\end{samepage}

\begin{proof}
A missing reachable consequence makes some row divergence infinite,
contradicting Proposition~\ref{prop:optimal-choice}.  A private consequence
can receive positive probability only if its action does, proving
\textup{(i)}.  Substituting $q^*=\delta_a$ into that proposition gives
\textup{(ii)}.  Under strict inequalities, any $r\ne\delta_a$ has
$\sum_b r(b)\Phi_b(K_a)<\bar u(a)$; \eqref{eq:choice-reference} then
gives strict suboptimality.  A rival outside $\supp K_a$ makes its
divergence infinite and violates the necessary inequality.
\end{proof}

\paragraph{Nonpositive trace weights.}
If $\lambda<0$, then
$V(q,K)\le\sum_aq(a)\bar u(a)\le\max_a\bar u(a)$.
Equality holds exactly when $q$ is supported on payoff-maximising actions
and $I_{qK}=0$.  Every payoff-maximising pure action attains that bound.
For $\lambda=0$, optimality requires only support on payoff maximisers.

\begin{proof}[Proof of Corollary~\ref{rem:logit}]
Disjoint row supports make $Z$ identify the action, so $I(A;Z)=\Sh(q)$.
Each action has a private consequence, hence the optimum has full support.
On $\supp K_a$, $p_{qK}=q(a)K_a$, so
$\KL(K_a\|p_{qK})=-\log q(a)$.  The optimality equalities become
$\bar u(a)-\lambda\log q^*(a)=c$.
Solving and normalising gives $q^*(a)\propto b_0^{\bar u(a)/\lambda}$;
in natural-log units this is the stated exponential formula.
Strict concavity of $\lambda\Sh(q)+\sum_aq(a)\bar u(a)$ proves uniqueness.
Changing the logarithm base merely rescales the information coefficient.
\end{proof}

\begin{proposition}[Optimal marginal and multiplicity]\label{prop:optimal-marginal}
Write $C^*(K):=\max_qV(q,K)$.  All optimal lotteries induce one output
distribution $p^*$.  With
\[
\mathcal A^*:=\{a:\bar u(a)+\lambda\KL(K_a\|p^*)=C^*(K)\},
\]
their set is exactly
\[
\{q\in\D(A):p_{qK}=p^*,\ \supp q\subseteq\mathcal A^*\}.
\]
The optimum is unique if the rows $\{K_a:a\in\mathcal A^*\}$ are affinely
independent.  Within any class of identical rows, value depends only on
the class's total probability.  The optimal lotteries are exactly the
lifts of the optimal lotteries of the collapsed channel, with each class's
mass split arbitrarily among its members.
\end{proposition}

\begin{proof}
Take one optimum with output $p^*$, which is positive on every row support.
Optimality gives $\Phi_a(p^*)\le C^*(K)$ for every $a$.  By
\eqref{eq:choice-reference}, a lottery $q$ attains $C^*(K)$ exactly when
both inequalities
\[
V(q,K)\le\sum_aq(a)\Phi_a(p^*)\le C^*(K)
\]
are equalities.  The first requires $p_{qK}=p^*$, and the second requires
$\supp q\subseteq\mathcal A^*$.  This proves output uniqueness and the
claimed optimal set.  Affine independence makes the mixture coefficients
representing $p^*$ unique.  Finally, identical rows have identical
$\bar u(a)$ and $\Sh(K_a)$, so \eqref{eq:choice-value} depends only on
their aggregate masses.  Collapsing or splitting a class therefore
preserves both value and optimality.
If all rows coincide, all lotteries consequently have the same value.
\end{proof}

\begin{remark}[Overlapping traces and IIA]\label{rem:overlap-iia-example}
Let three actions have equal expected utility and rows
\[
K_1=(1,0,0),\qquad K_2=(0,1,0),\qquad K_3=(1/2,0,1/2).
\]
With the first two actions only, the optimum is $(1/2,1/2)$.  Adding the
third gives
\[
q^*=(1/3,4/9,2/9),\qquad p_{q^*K}=(4/9,4/9,1/9).
\]
All three row divergences equal $\log(9/4)$.
Proposition~\ref{prop:optimal-choice} verifies optimality; the rows are
affinely independent, so it is unique.  The first-to-second choice ratio
falls from $1$ to $3/4$, exhibiting IIA failure with equal expected utilities.
\end{remark}

\subsubsection{The value of being on record}
\label{app:record-value-proofs}

\begin{proposition}[Cost of a record garbling]\label{prop:garbling-price}
Let $L=KT$ for a payoff-preserving record processor $T$, and fix $q$.
Under the joint law $q(a)K(o,r\mid a)T(o,r'\mid o,r)$,
\[
V(q,K)-V(q,L)=\lambda I(A;R\mid O,R')\ge0.
\]
Equality holds exactly when, after $(O,R')$ is observed, $R$ contains
no further information about the action.
\end{proposition}

\begin{proof}
The payoff terms agree.  Put $Z=(O,R)$ and $Z'=(O,R')$.
The Markov chain $A-Z-Z'$ and the two chain-rule expansions of
$I(A;Z,Z')$ give
\[
I(A;Z)-I(A;Z')=I(A;Z\mid Z')=I(A;R\mid O,R').
\]
Conditional mutual information is the average divergence between the
conditional joint distribution and the product of its marginals.  It is
therefore nonnegative, and vanishes exactly at the stated conditional
independence.  Under this independence the Markov property gives
$\pi^{qK}_Z=\Pr(A=\cdot\mid Z,Z')=\pi^{qL}_{Z'}$ almost surely,
so the posterior laws also agree.  Conversely, equality of those laws
gives equal information through
$I(A;Z)=\Sh(q)-\E[\Sh(\pi^{qK}_Z)]$.
Thus zero garbling cost is also equivalent to $\mu_{qK}=\mu_{qL}$.
\end{proof}

\begin{corollary}[The recorded coin]\label{cor:recorded-coin}
For $K,L:A\to\D(O\times R)$ and $t\in(0,1)$, let
$tK\oplus(1-t)L$ draw an action-independent coin $Y$ with
$\Pr(Y=0)=t$, run $K$ or $L$ accordingly, and retain $(Y,R)$ in the
record.  Let $M_t:=tK+(1-t)L$ hide that coin.  For every $q$,
\begin{enumerate}[label=\textup{(\roman*)},leftmargin=*]
\item $V(q,tK\oplus(1-t)L)=tV(q,K)+(1-t)V(q,L)$;
\item $V(q,tK\oplus(1-t)L)-V(q,M_t)=\lambda I(A;Y\mid O,R)\ge0$,
with equality exactly when $A$ and $Y$ are conditionally independent
given $(O,R)$.
\end{enumerate}
\end{corollary}

\begin{proof}
For \textup{(i)}, expected utility averages across the coin and
$I(A;Y)=0$, so the chain rule averages the two information terms as well.
For \textup{(ii)}, deleting $Y$ is the record garbling from the public
mixture to $M_t$; apply Proposition~\ref{prop:garbling-price}.
\end{proof}

\begin{samepage}
\begin{proposition}[Uniform dominance is more-capable, not garbling]
\label{prop:blackwell}\label{def:channel-orders}
Let $K:A\to\D(O\times R)$ and $L:A\to\D(O\times R')$ have the same
payoff kernels.  Write $K\succeq_gL$ if $L=KT$ for some record processor,
and $K\succeq_cL$ if $I_{qK}(A;O,R)\ge I_{qL}(A;O,R')$ for every $q$.
The latter is the \emph{more-capable} order \cite{KornerMarton1977}.
\begin{enumerate}[label=\textup{(\roman*)},leftmargin=*]
\item $K\succeq_cL$ exactly when $V(q,K)\ge V(q,L)$ for every $q$.
\item $K\succeq_gL$ implies $K\succeq_cL$, but the converse fails:
there are such channels with $K\succeq_cL$ and no record processor
carrying $K$ to $L$.
\end{enumerate}
\end{proposition}
\end{samepage}

\begin{proof}
Identical payoff kernels give
$V(q,K)-V(q,L)=\lambda[I_{qK}-I_{qL}]$, proving \textup{(i)}.
The garbling implication is Proposition~\ref{prop:garbling-price}.
For the converse, let $A=\{0,1\}$ and fix one payoff $\bar o$ under
both channels.  Let $K$ reveal the action with probability $1/2$ and
otherwise report an erasure $\mathrm e$; let $L$ report the action with
probability $4/5$ and its opposite with probability $1/5$.
In bits, with binary entropy $h$ and $q=(1-t,t)$,
\[
I_{qK}=\tfrac12h(t),\qquad
I_{qL}=h\bigl(\tfrac15+\tfrac35t\bigr)-h(\tfrac15).
\]
Put $d(t):=I_{qK}-I_{qL}$.  It is continuous on $[0,1]$, with
$d(0)=d(1)=0$, and elementary differentiation gives
\[
d''(t)=\frac{9t(1-t)-4}
 {2\ln2\,t(1-t)[4+9t(1-t)]}<0\qquad(0<t<1),
\]
because $t(1-t)\le1/4$.  Concavity therefore gives $d(t)\ge0$,
proving $K\succeq_cL$ in every logarithm base.
Suppress the fixed payoff coordinate.  If a processor $T$ carried $K$ to
$L$, subtracting its output-$0$
probabilities under the two actions would give
\[
\tfrac35=\tfrac12\bigl[T(0\mid0)-T(0\mid1)\bigr],
\]
which is impossible since the right-hand side is at most $1/2$.
The erasure contribution cancels because it is the same in both rows.
Compositions of processors are processors, so no chain of garblings
can produce $L$ from $K$ either.
\end{proof}

\subsubsection{Measuring the trace weight}
\label{app:measurement-proofs}

\begin{proposition}[Single information price]
\label{prop:single-price}\label{lem:info-payoff-plane}
If $(q,K)\sim(p,L)$ and $\I_{qK}\ne\I_{pL}$, then
\[
 \frac{E_{pL}-E_{qK}}{\I_{qK}-\I_{pL}}=\lambda.
\]
Only $\lambda$ satisfies this identity at every such indifference.
More generally, put $\underline u:=\min_o u(o)$ and
$\overline u:=\max_o u(o)$. The attainable pairs
$(\I_{qK},E_{qK})$, over all finite experiments, form exactly
$[0,\infty)\times[\underline u,\overline u]$. Indifference classes are
the intersections of this strip with the lines $E=c-\lambda\I$.
\end{proposition}

\begin{proof}
The ratio rearranges $E_{qK}+\lambda\I_{qK}=E_{pL}+\lambda\I_{pL}$.
For attainability, fix $(t,w)$ in the strip and choose $n\ge2$ with
$\log n\ge t$. Under the uniform action lottery, reveal the action with
probability $t/\log n$ and otherwise report an erasure. Independently pay
a lottery over utility-minimising and utility-maximising outcomes with
expected utility $w$. This gives information $t$ and material value $w$.
The representation now gives the stated indifference lines.
Finally, choose $0<t<(\overline u-\underline u)/\lambda$.
The attainable pairs $(t,\overline u-\lambda t)$ and $(0,\overline u)$
are indifferent and have different information, so their ratio identifies
$\lambda$ uniquely.
\end{proof}

\begin{samepage}
\begin{proposition}[Willingness to pay for trace]
\label{prop:wtp}\label{def:dilution}
Choose payoffs $o^+,o^-$ with $u(o^+)>u(o^-)$ and put
\[
 \Delta u:=u(o^+)-u(o^-)>0,\qquad
 w(\beta):=\beta u(o^+)+(1-\beta)u(o^-).
\]
Let $J_\beta$ pay $o^+$ with probability $\beta$ and $o^-$ otherwise,
independently of the action, with a constant record. For any channel $N$,
let $D_{\varepsilon,\beta}N$ be the public mixture that runs $N$ with
probability $1-\varepsilon$ and $J_\beta$ with probability $\varepsilon$,
retaining the branch label in the record. Here $0\le\varepsilon<1$ and
$0\le\beta\le1$; the two channels use the same action set.
For any experiments $(q,K)$ and $(p,L)$:
\begin{enumerate}[label=\textup{(\roman*)},leftmargin=*]
\item
\[
 V(q,D_{\varepsilon,\beta}K)
 =(1-\varepsilon)V(q,K)+\varepsilon w(\beta).
\]
\item If $\varepsilon>0$, indifference between
$(q,D_{\varepsilon,\beta}K)$ and $(p,D_{\varepsilon,\beta'}L)$ is equivalent to
\[
 \frac{\varepsilon}{1-\varepsilon}
 [w(\beta)-w(\beta')]=V(p,L)-V(q,K).
\]
If $E_{qK}=E_{pL}$, the right side is
$\lambda(\I_{pL}-\I_{qK})$.
\item Such an indifference can be obtained with $\beta,\beta'\in[0,1]$
if and only if
\[
 |V(p,L)-V(q,K)|\le
 \frac{\varepsilon}{1-\varepsilon}\Delta u.
\]
Strict preference even at maximal compensation gives the corresponding
strict bound on the value difference.
\end{enumerate}
\end{proposition}
\end{samepage}

\begin{proof}
The case $\varepsilon=0$ is immediate. For $\varepsilon>0$,
the public-mixture identity in Corollary~\ref{cor:recorded-coin} proves
\textup{(i)}, since $J_\beta$ has value $w(\beta)$ and no information.
Subtract the two values to obtain \textup{(ii)}.
Their compensation term ranges over
$[-\varepsilon\Delta u/(1-\varepsilon),
\varepsilon\Delta u/(1-\varepsilon)]$, proving \textup{(iii)}.
\end{proof}

\begin{proof}[Proof of Proposition~\ref{prop:elicit}]
Write $K^{\mathrm{id}}$ and $K^0$ for revelation and silence at the common
payoff. Their value difference is $\lambda\Sh(q)$.
Proposition~\ref{prop:wtp}, with $u(o^-)=0$ and $u(o^+)=1$, makes the
diluted revealing-minus-silent difference
\[
 (1-\varepsilon)\lambda\Sh(q)-\varepsilon\beta.
\]
It has a unique zero in $[0,1]$ exactly when
$\lambda\Sh(q)\le\varepsilon/(1-\varepsilon)$; otherwise it is positive
throughout. Solving for its zero gives the elicitation formula.
\end{proof}

\paragraph{The two-prior test.}
At a common $\varepsilon$, two indifference reports satisfy
$\beta_1/\beta_2=\Sh(q_1)/\Sh(q_2)$.
For $q_1=(1/2,1/2)$ and $q_2=(1/4,3/4)$ this is
$\log2/[\frac14\log4+\frac34\log(4/3)]\approx1.23$.
Under the Gini-reduction criterion of Appendix~\ref{app:axiom-independence},
the same public mixture weights information value by $1-\varepsilon$,
but revelation is worth $G(q)=1-\sum_aq(a)^2$.
It therefore predicts $(1/2)/(3/8)=4/3$ instead.

\begin{remark}[Implementation and controls]\label{rem:elicitation}
Choose utility-extreme payoffs and normalise $u(o^-)=0$, $u(o^+)=1$.
To measure $u(o)$, compare the sure payoff $o$ with $J_\beta$ on a
single-action set: their indifference gives $u(o)=\beta$.
For uniform $q$ on $n\ge2$ actions, elicit
\[
 (q,D_{\varepsilon,\beta}K_o^0)
 \sim(q,D_{\varepsilon,0}K_o^{\mathrm{id}}),
 \qquad
 \lambda=\frac{\varepsilon\beta}{(1-\varepsilon)\log n},
\]
where the two channels pay $o$ surely and respectively hide and reveal
the action. If revelation is preferred even at $\beta=1$, then
$\lambda>\varepsilon/[(1-\varepsilon)\log n]$; increasing $\varepsilon$
eventually permits indifference.

Three restrictions provide controls. Revelation must beat silence at
every constant payoff. Erasing the records of these two undiluted channels
must make them indifferent. Finally, every successful elicitation must
return the same $\lambda$ across priors, baseline payoffs, and action-set
sizes. The capped criterion $E+\min\{\I,c\}$ instead returns
$\min\{1,c/[(1-\varepsilon)\Sh(q)]\}$ whenever compensation is feasible.
Thus a falling estimate at high entropy is consistent with satiation,
and rejects a single information price.
\end{remark}

\subsubsection{Moral hazard and non-contractible records}
\label{app:moral-hazard-proofs}

\paragraph{The effort problem.}
Let $K_a$ be Bernoulli with success probability $p_a$, where
$0<p_0<p_1<1$, put $d=p_1-p_0$, and let $q$ be the probability of high
effort.  Write $h(t)=-t\log t-(1-t)\log(1-t)$ for binary entropy, with
the paper's fixed logarithm base $b>1$.  Then
\[
 I(q):=\I_{qK}(A;X)
 =h(p_0+qd)-(1-q)h(p_0)-qh(p_1),
\]
and the agent's value is
$V(q)=w_0+(p_0+qd)\Delta w-qc+\lambda I(q)$.
For $0<q<1$,
\[
 I''(q)=-\frac{d^2}{\ln b\,(p_0+qd)(1-p_0-qd)}<0,
 \qquad
 I'(0)=\KL(K_1\|K_0),\quad
 I'(1)=-\KL(K_0\|K_1).
\]
Thus $V$ is strictly concave.  The endpoint derivative at one is
$d\Delta w-c-\lambda\KL(K_0\|K_1)$, so pure high effort is optimal
exactly when
\[
 \Delta w\ge\frac{c+\lambda\KL(K_0\|K_1)}d.
\]
At a flat wage, $V'(1)<0$ and
$V'(0)=-c+\lambda\KL(K_1\|K_0)$; hence the unique optimum is interior
exactly when $c<\lambda\KL(K_1\|K_0)$, and otherwise it is zero.
These calculations apply even when the high-effort action has lower
expected material utility, without relabelling the actions.

The cost implementation in the main text preserves these comparisons.
For each candidate plan $q$, independently draw $Z_q\in\{0,c\}$ with
$\Pr(Z_q=c)=q$ and use material utility $u(w,z)=w-z$.
This gives expected cost $qc$.  Since $W=w_0+\Delta wX$ is a function
of $X$ and $Z_q$ is independent of $(A,X)$,
\[
 \I(A;W,X,Z_q)=\I(A;X)+\I(A;Z_q\mid X)=\I(A;X).
\]
The augmented channel varies with the plan; its value therefore follows
from the theorem's cross-channel comparison clause, rather than from
optimisation within one fixed augmented channel.

\paragraph{Teams.}
Hold the teammates' joint distribution fixed and independent of the
agent's own randomisation.  Conditional on her action $a$, integrate their
efforts out of $g(a,A_{-i})$.  This produces a Bernoulli row with parameter
$\bar g_a=\E[g(a,A_{-i})]$.  If $\bar g_1>\bar g_0$, the preceding
calculation applies verbatim.  The uniform bounds on $g$ keep these
parameters away from zero and one; positive marginal productivity is a
separate assumption.

For a sequence of such team problems, put $p=\bar g_0$ and
$d=\bar g_1-\bar g_0>0$.  Uniformly when the probabilities remain in
$[\varepsilon,1-\varepsilon]$, Taylor's theorem gives
\[
 \KL\bigl(\operatorname{Bern}(p)\|\operatorname{Bern}(p+d)\bigr)
 =\frac{d^2}{2\ln b\,p(1-p)}+O(d^3).
\]
Indeed, the divergence is zero at $d=0$, its first derivative is zero
there, and its second derivative is $1/(\ln b\,p(1-p))$; its third
derivative is uniformly bounded on the stated probability range.
Consequently the required bonus has cost component $c/d$ and trace
component
\[
 \frac{\lambda\KL(\bar K_0\|\bar K_1)}d
 =\frac{\lambda d}{2\ln b\,p(1-p)}+O(d^2),
\]
so the latter vanishes as $d\to0$.  For example, if
$g_n(a)=G(n^{-1}\sum_i a_i)$ and $G$ is increasing and Lipschitz with
constant $L$ independent of $n$, then $0\le d\le L/n$.
Dependence on the fraction working without such regularity does not by
itself imply a vanishing marginal effect: a threshold technology can have
a jump.  Nor does the conclusion apply to an essential task whose marginal
effect stays bounded away from zero.

\paragraph{Adding a signal and adjusting the wage.}
Let $L(x,s\mid a)$ be the joint output--signal channel, whose $X$ marginal
is the original production channel.  Wages still depend only on $X$.
For each fixed $q$, the chain rule gives the additional utility
\[
 G(q):=\lambda\I_q(A;S\mid X)\ge0.
\]
To allow effort to adjust, let
$J(q)=(p_0+qd)\Delta w-qc+\lambda\I_q(A;X)$ be the old value excluding
the base wage, and define
\[
 \Gamma:=\max_{q\in[0,1]}[J(q)+G(q)]-\max_{q\in[0,1]}J(q).
\]
If $q^*$ maximises $J$, then
\[
 G(q^*)\le\Gamma\le\max_q G(q).
\]
The lower bound evaluates the new objective at $q^*$; the upper bound
uses $J(q)\le\max J$.  A reduction of the base wage by $\Gamma$ therefore
preserves the agent's maximised utility exactly.  Reducing it by
$G(q^*)$ preserves participation weakly, assuming the wage reduction is
feasible.  A base-wage change does not affect the optimal effort lottery,
but the added signal generally does.  Thus this argument establishes
compensation and participation, not preservation of effort or an
unconditional increase in the principal's profit.  When $\lambda=0$,
$G$ and $\Gamma$ are zero.

\subsection{Interpretation: prediction, control, and an audience}
\label{app:interpretation-proofs}

\paragraph{What information measures.}
Write $Z=(O,R)$, $K_a=K(\cdot\mid a)$, and $p(z)=\sum_aq(a)K_a(z)$.
With $q_z=\Pr(A=\cdot\mid Z=z)$ on reached outputs and zero-weight
terms omitted, the finite identities
\[
 \I_q(A;Z)
 =\Sh(q)-\sum_zp(z)\Sh(q_z)
 =\Sh(p)-\sum_aq(a)\Sh(K_a)
 =\sum_aq(a)\KL(K_a\|p)
\]
justify the information claims in the introduction and
Section~\ref{sec:interpretation}.  They follow by expanding the logarithm
in \eqref{eq:mi}.  Nonnegativity of divergence shows that information is
zero exactly when the supported channel rows coincide, equivalently when
$A$ and $Z$ are independent.  A pure action therefore has zero trace.
Pairwise-disjoint supported rows make every reached posterior a point mass,
giving information $\Sh(q)$.  The same joint-law definition is symmetric
in $A$ and $Z$, so reading the dependence as evidence about the action or
as the action's influence on consequences gives the same index and rankings.

If all actions have the same expected payoff utility $\bar u$, then
$V(q,K)=\bar u+\lambda\I_q(A;Z)$.  Since $\lambda>0$, its maximisers are
exactly the distributions attaining channel capacity
$C(K)=\max_q\I_q(A;Z)$, and the maximal value is
$\bar u+\lambda C(K)$.  This proves the stated relation to one-step
empowerment without identifying capacity with the value of an arbitrary
fixed lottery.

\paragraph{Prediction error.}
For any forecast distribution $r\in\D(Z)$, the log-score identity is
\[
 \sum_aq(a)\KL(K_a\|r)
 =\I_q(A;Z)+\KL(p\|r).
\]
Terms with positive mass and zero forecast probability are infinite.
Minimisation over $r$ therefore gives $\I_q(A;Z)$, attained at $r=p$.
Equivalently, minimum expected log loss without observing the action is
$\Sh(p)$, whereas knowing it permits the row-specific forecast $K_a$ and
reduces this loss to $\sum_aq(a)\Sh(K_a)$ on average.
Trace is precisely this \emph{excess} prediction error from not knowing
the action.  It is not unadjusted outcome uncertainty: identical noisy
rows can have arbitrarily high outcome entropy and zero trace.
Within a fixed channel, the criterion can thus be written as expected
material utility plus the minimum expected forecasting loss
$\lambda\KL(K_a\|r)$.  This is the precise forecasting interpretation used
in Section~\ref{sec:literature}.

Any reclassification of payoff and record that preserves the complete
visible information leaves mutual information unchanged: the old and new
visible variables are obtainable from one another, so data processing in
both directions gives equal mutual information.  Likewise, adding a passive
audience while keeping $q$, $K$, and the material payoffs fixed changes no
argument of $V$.

\begin{proposition}[An entropy-reduction audience]\label{prop:audience}
Fix an actual action lottery $q$ and an observer's prior $r$ on the same
finite action set, with $\supp q\subseteq\supp r$.  The observer knows
the channel $K:A\to\D(Z)$ and, after output $z$, forms its Bayesian
posterior $r_z$.  Let its reward, in utility units, be
$\lambda[\Sh(r)-\Sh(r_z)]$, where $\lambda\ne0$.
Its actual expected reward equals $\lambda\I_{qK}(A;Z)$ for every finite
channel $K$ if and only if $r=q$.
\end{proposition}

\begin{proof}
The support condition ensures that the observer's posterior is defined
at every output with positive actual probability.  The expected reward is
\[
 \lambda\left[\Sh(r)-\sum_zp_{qK}(z)\Sh(r_z)\right].
\]
For $r=q$, the posterior-entropy identity proves the claim.

Conversely, suppose the equality holds for every channel and cancel
$\lambda\ne0$.  Full revelation gives $\Sh(r)=\Sh(q)$.
Choose $b$ with $q(b)>0$, so that $r(b)>0$, and set
$x=\tfrac12\min\{q(b),r(b)\}\in(0,1/2]$.
For each $a\ne b$, consider the channel that always reports a new record
$\ast$ after $b$, reports $\ast$ with probability $x$ after $a$, and
otherwise reveals the action.
For either prior $p\in\{q,r\}$, only $\ast$ can leave uncertainty; the
posterior probability of $a$ at that record is
\[
 t_p:=\frac{xp(a)}{xp(a)+p(b)}.
\]
Let $h(t):=\Sh(t,1-t)$ denote binary entropy.  The actual probability of
$\ast$ is $xq(a)+q(b)>0$.  Since the prior entropies agree, the assumed
reward identity for this channel gives $h(t_r)=h(t_q)$.
Both posterior probabilities lie in $[0,1/2]$, because
$xp(a)\le x\le p(b)$ for $p=q,r$.  Binary entropy is strictly increasing
on that interval, so $t_r=t_q$, or
\[
 q(a)r(b)=r(a)q(b).
\]
This identity also holds for $a=b$.  Summing over $a$ gives $r(b)=q(b)$;
since $q(b)>0$, the same identity then gives $r(a)=q(a)$ for every $a$.
\end{proof}

To reproduce the criterion on a whole menu of plans, the observer must
therefore use each implemented lottery as its prior, for example because
the lottery is observed or announced.  Correct anticipation at a single
equilibrium lottery gives the equality there; it does not establish the
same formula for deviations evaluated against a fixed observer prior.
The rewards above are additive utility rewards.  They can be monetary
payments when the agent is quasilinear in those payments; expected money
alone does not give the stated utility formula under general payoff risk
preferences.

\clearpage
\begin{thebibliography}{99}
\raggedright

\bibitem{AgranovOrtoleva2017}
Marina Agranov and Pietro Ortoleva.
\newblock Stochastic choice and preferences for randomization.
\newblock {\em Journal of Political Economy}, 125(1):40--68, 2017.

\bibitem{AndreoniBernheim2009}
James Andreoni and B. Douglas Bernheim.
\newblock Social image and the 50--50 norm: A theoretical and experimental analysis of audience effects.
\newblock {\em Econometrica}, 77(5):1607--1636, 2009.

\bibitem{AzrieliLehrer2008}
Yaron Azrieli and Ehud Lehrer.
\newblock The value of a stochastic information structure.
\newblock {\em Games and Economic Behavior}, 63(2):679--693, 2008.

\bibitem{BartlingFehrHerz2014}
Bj\"orn Bartling, Ernst Fehr, and Holger Herz.
\newblock The intrinsic value of decision rights.
\newblock {\em Econometrica}, 82(6):2005--2039, 2014.

\bibitem{BaezFritzLeinster2011}
John C. Baez, Tobias Fritz, and Tom Leinster.
\newblock A characterization of entropy in terms of information loss.
\newblock {\em Entropy}, 13(11):1945--1957, 2011.

\bibitem{BenabouTirole2002}
Roland B\'enabou and Jean Tirole.
\newblock Self-confidence and personal motivation.
\newblock {\em Quarterly Journal of Economics}, 117(3):871--915, 2002.

\bibitem{BenabouTirole2006}
Roland B\'enabou and Jean Tirole.
\newblock Incentives and prosocial behavior.
\newblock {\em American Economic Review}, 96(5):1652--1678, 2006.

\bibitem{Bernheim1994}
B. Douglas Bernheim.
\newblock A theory of conformity.
\newblock {\em Journal of Political Economy}, 102(5):841--877, 1994.

\bibitem{Blackwell1953}
David Blackwell.
\newblock Equivalent comparisons of experiments.
\newblock {\em Annals of Mathematical Statistics}, 24(2):265--272, 1953.


\bibitem{CabralesGossnerSerrano2013}
Antonio Cabrales, Olivier Gossner, and Roberto Serrano.
\newblock Entropy and the value of information for investors.
\newblock {\em American Economic Review}, 103(1):360--377, 2013.

\bibitem{CaplinDeanLeahy2022}
Andrew Caplin, Mark Dean, and John Leahy.
\newblock Rationally inattentive behavior: Characterizing and generalizing {S}hannon entropy.
\newblock {\em Journal of Political Economy}, 130(6):1676--1715, 2022.

\bibitem{CaplinLeahy2001}
Andrew Caplin and John Leahy.
\newblock Psychological expected utility theory and anticipatory feelings.
\newblock {\em Quarterly Journal of Economics}, 116(1):55--79, 2001.

\bibitem{CerreiaVioglioDillenbergerOrtolevaRiella2019}
Simone Cerreia-Vioglio, David Dillenberger, Pietro Ortoleva, and Gil Riella.
\newblock Deliberately stochastic.
\newblock {\em American Economic Review}, 109(7):2425--2445, 2019.

\bibitem{CorraoFudenbergMudekerezaLevine2026}
Roberto Corrao, Drew Fudenberg, Florian Mudekereza, and David Levine.
\newblock Randomization, surprise, and adversarial forecasters.
\newblock Working paper, August 2026.

\bibitem{CoverThomas2006}
Thomas M. Cover and Joy A. Thomas.
\newblock {\em Elements of Information Theory}.
\newblock Wiley-Interscience, Hoboken, second edition, 2006.

\bibitem{Denti2022}
Tommaso Denti.
\newblock Posterior separable cost of information.
\newblock {\em American Economic Review}, 112(10):3215--3259, 2022.

\bibitem{Dostoevsky1864}
Fyodor Dostoevsky.
\newblock {\em Notes from the Underground}, Part I.
\newblock 1864.  Translated by Constance Garnett, 1918.

\bibitem{ElyFrankelKamenica2015}
Jeffrey Ely, Alexander Frankel, and Emir Kamenica.
\newblock Suspense and surprise.
\newblock {\em Journal of Political Economy}, 123(1):215--260, 2015.

\bibitem{Faddeev1956}
D. K. Faddeev.
\newblock On the concept of entropy of a finite probabilistic scheme.
\newblock {\em Uspekhi Matematicheskikh Nauk}, 11(1):227--231, 1956.
\newblock In Russian.

\bibitem{FrankelKamenica2019}
Alexander Frankel and Emir Kamenica.
\newblock Quantifying information and uncertainty.
\newblock {\em American Economic Review}, 109(10):3650--3680, 2019.

\bibitem{GilboaLehrer1991}
Itzhak Gilboa and Ehud Lehrer.
\newblock The value of information---an axiomatic approach.
\newblock {\em Journal of Mathematical Economics}, 20(5):443--459, 1991.

\bibitem{GrantKajiiPolak1998}
Simon Grant, Atsushi Kajii, and Ben Polak.
\newblock Intrinsic preference for information.
\newblock {\em Journal of Economic Theory}, 83(2):233--259, 1998.

\bibitem{HersteinMilnor1953}
I. N. Herstein and John Milnor.
\newblock An axiomatic approach to measurable utility.
\newblock {\em Econometrica}, 21(2):291--297, 1953.

\bibitem{Holmstrom1979}
Bengt Holmstr\"om.
\newblock Moral hazard and observability.
\newblock {\em Bell Journal of Economics}, 10(1):74--91, 1979.

\bibitem{Holmstrom1982}
Bengt Holmstr\"om.
\newblock Moral hazard in teams.
\newblock {\em Bell Journal of Economics}, 13(2):324--340, 1982.


\bibitem{KlyubinPolaniNehaniv2005}
Alexander S. Klyubin, Daniel Polani, and Chrystopher L. Nehaniv.
\newblock Empowerment: A universal agent-centric measure of control.
\newblock In {\em Proceedings of the 2005 IEEE Congress on Evolutionary Computation}, volume~1, pages 128--135. IEEE,
  2005.

\bibitem{KornerMarton1977}
J\'anos K\"orner and Katalin Marton.
\newblock Comparison of two noisy channels.
\newblock In Imre Csisz\'ar and Peter Elias, editors, {\em Topics in Information Theory}, Colloquia Mathematica
  Societatis J\'anos Bolyai, volume~16, pages 411--423. North-Holland, Amsterdam, 1977.


\bibitem{MatejkaMcKay2015}
Filip Mat\v{e}jka and Alisdair McKay.
\newblock Rational inattention to discrete choices: A new foundation for the multinomial logit model.
\newblock {\em American Economic Review}, 105(1):272--298, 2015.

\bibitem{MistryLiljeholm2016}
Prachi Mistry and Mimi Liljeholm.
\newblock Instrumental divergence and the value of control.
\newblock {\em Scientific Reports}, 6:36295, 2016.

\bibitem{NortonLiljeholm2020}
Kaitlyn G. Norton and Mimi Liljeholm.
\newblock The rostrolateral prefrontal cortex mediates a preference for high-agency environments.
\newblock {\em Journal of Neuroscience}, 40(22):4401--4409, 2020.


\bibitem{PomattoStrackTamuz2023}
Luciano Pomatto, Philipp Strack, and Omer Tamuz.
\newblock The cost of information: The case of constant marginal costs.
\newblock {\em American Economic Review}, 113(5):1360--1393, 2023.

\bibitem{SalgeGlackinPolani2014}
Christoph Salge, Cornelius Glackin, and Daniel Polani.
\newblock Empowerment---an introduction.
\newblock In Mikhail Prokopenko, editor, {\em Guided Self-Organization: Inception}, pages 67--114. Springer, Berlin,
  2014.

\bibitem{Savage1954}
Leonard J. Savage.
\newblock {\em The Foundations of Statistics}.
\newblock Wiley, New York, 1954.

\bibitem{Shannon1948}
Claude E. Shannon.
\newblock A mathematical theory of communication.
\newblock {\em Bell System Technical Journal}, 27:379--423, 623--656, 1948.

\bibitem{Sims2003}
Christopher A. Sims.
\newblock Implications of rational inattention.
\newblock {\em Journal of Monetary Economics}, 50(3):665--690, 2003.

\bibitem{WynerZiv1973}
Aaron D. Wyner and Jacob Ziv.
\newblock A theorem on the entropy of certain binary sequences and applications: Part {I}.
\newblock {\em IEEE Transactions on Information Theory}, 19(6):769--772, 1973.

\end{thebibliography}


\end{document}
